\chapter{The index theorem via semiclassical and adiabatic limits}\label{ch:semiclassical-index}
\chaptermark{The semiclassical index}
\providecommand{\Kc}{K_{\mathrm{c}}}
\providecommand{\inda}{\operatorname{ind}_{\mathrm{a}}}
\providecommand{\indt}{\operatorname{ind}_{\mathrm{t}}}
\providecommand{\Ginf}{G^{-\infty}}
\providecommand{\Ch}{\operatorname{Ch}}
\providecommand{\Td}{\operatorname{Td}}
\providecommand{\Gs}[1]{\Gamma^{#1}}
\providecommand{\Indsl}[1]{\operatorname{Ind}^{#1}_{\mathrm{sl}}}
\providecommand{\indT}{\operatorname{ind}_{\mathcal T}}
\providecommand{\Psisl}[1]{\Psi^{#1}_{\mathrm{sl}}}
\providecommand{\Psiinf}{\Psi^{-\infty}_{\epsilon^\infty}}
\providecommand{\sigsl}{\sigma_{\mathrm{sl}}}
\providecommand{\Ope}{\operatorname{Op}_\epsilon}
\providecommand{\Sc}{\mathcal{S}_{\mathrm c}}
\providecommand{\Torus}{\mathbb{T}}
\providecommand{\Hev}{H^{\mathrm{ev}}}
\providecommand{\Ssl}[1]{S^{#1}_{\mathrm{sl}}}

In \cref{ch:families} the families index theorem of Atiyah and Singer was proved
in K-theory by embedding a fibration $\phi:M\to B$ into a trivial fibration
$\R^N\times B\to B$ and using three properties of the analytic index:
multiplicativity for towers of fibrations, excision, and a normalization on the
trivial fibration. One input, the index of the Thom class on a sphere bundle
(statement \ref{families:TH} there), was imported. This chapter gives a second proof, based on
a different realization of the index map: the \emph{semiclassical index}.

The idea is simple. A compactly supported K-class on $T^*(M/B)$ can be
represented by data which are trivial outside a compact set; for idempotents,
say, by a smooth family $e$ of idempotent matrices which is equal to a constant
idempotent $p_0$ outside a compact subset of $T^*(M/B)$. Such a symbol can be
quantized, with a small parameter $\epsilon>0$ (Planck's constant), to a
family of idempotent operators on the fibres which differs from $p_0$ by a
\emph{semiclassical family of smoothing operators}. A pair of idempotents
differing by a smoothing operator has a relative index, an element of $K^0(B)$,
and this is the semiclassical index $\Indsl{}$. As $\epsilon\downarrow0$ the
Schwartz kernels concentrate at the fibre diagonal, so the construction is
local: excision for $\Indsl{}$ is immediate. The main analytic result of the
chapter (\cref{scindex:thm:analytic=sl}) is that $\Indsl{}$ \emph{coincides} with
the analytic index of \cref{ch:product,ch:families}. The proof is short and
algebraic: from an elliptic family and a parametrix one forms an explicit
idempotent (\eqref{scindex:eq:milnor}) whose relative index is the index bundle;
if the family and the parametrix are themselves semiclassical quantizations, this
idempotent is a semiclassical quantization of the corresponding symbol-level
idempotent.

The second ingredient is the index of the Thom class. In the semiclassical
picture this becomes a statement about the harmonic oscillator on the fibres of
the normal bundle: after rescaling by $\epsilon^{1/2}$, the semiclassical
quantization of the Thom symbol is the Koszul complex of the harmonic
oscillator, whose null space is spanned by the Gaussian. Since the Gaussian is
invariant under the orthogonal group, its span is a \emph{canonically trivial}
line bundle; no equivariant index computation is needed
(\cref{scindex:thm:thom}). As a consequence, statement (TH) of
\cref{ch:families} (\ref{families:TH}) is proved, and the semiclassical normalization on the trivial
fibration follows as well.

Putting these together (\cref{scindex:sec:AS}) gives the families index theorem
in the following form: excision and the normalization are replaced by their
(essentially trivial) semiclassical versions together with
\cref{scindex:thm:analytic=sl}, and (TH) by \cref{scindex:thm:thom}; the only
property of the analytic index from \cref{ch:families} which is still used is
multiplicativity (\cref{families:thm:mult,families:cor:mult}). In R.~Melrose's treatment, on which this chapter is based, the
multiplicativity of the semiclassical index is also proved semiclassically, by an
adiabatic limit; we describe that argument in
\cref{scindex:sec:adiabatic}, without giving the full details of the
underlying two-parameter calculus.

The odd version of the semiclassical index, $\Kc^{-1}(T^*(M/B))\to K^{-1}(B)$,
obtained by quantizing compactly supported perturbations of the identity, is the
most elementary of these maps; we show in \cref{scindex:sec:odd} that the even
index is the odd one applied to loops, followed by the Toeplitz index
isomorphism of \cref{ch:bott} (\cref{bott:thm:index-iso}). Throughout, the K-groups of the base are realized,
as in \cref{ch:index,ch:bott}, through the group $\Ginf$ of invertible
perturbations of the identity by smoothing operators. The required passage from
families of operators on the varying fibres $Z_b$ to maps into a fixed group
$\Ginf(Y;F)$ is made in \cref{scindex:sec:smoothing} by an explicit isometric
embedding of the fibres into a torus.

\section{Setting and recollections}\label{scindex:sec:setting}

\subsection{Fibrations}
Throughout, $B$ is a compact manifold, possibly with boundary (we need
$B\times[0,1]$ for homotopies). A \emph{fibration over $B$} in this chapter is a
submersion $\phi:M\to B$ such that $M$ is an open subset of the total space of a
fibration $\bar\phi:\bar M\to B$ of compact manifolds in the sense of
\cref{ch:product} (\cref{product:def:fibration}), with $\phi=\bar\phi|_M$. The fibres $Z_b=\phi^{-1}(b)$ may
thus be non-compact. The cases we need are:
\begin{itemize}
\item fibrations of compact manifolds ($M=\bar M$);
\item open subsets $U\subset\bar M$;
\item the total space $\mathcal V$ of a real vector bundle $\varpi:\mathcal V\to
  M$ over the total space of a fibration of compact manifolds, fibred over $B$
  by $\phi\circ\varpi$; it is an open subset of the sphere bundle
  $S(\mathcal V\oplus\R)$, via inverse stereographic projection;
\item the trivial fibrations $\R^N\times B$ and $I^N\times B$, $I=(0,2\pi)$,
  inside $\Torus^N\times B$, $\Torus=\R/2\pi\Z$;
\item the degenerate fibration $B\to B$, $\phi=\Id$, whose fibres are points.
\end{itemize}
We write $M^2_\phi=M\times_BM$ for the fibre product, $\Diag\subset M^2_\phi$ for
the (fibre) diagonal, $T^*(M/B)$ for the vertical cotangent bundle and
$\varpi:T^*(M/B)\to M$ for its projection. A Hermitian bundle $E\to M$ and a
smooth positive fibre density $\nu$ are always fixed; they define the Hilbert
spaces $L^2(Z_b;E_b)$. We write $\Psido{-\infty}_{\mathrm c}(M/B;E)$ for the
smoothing families with compactly supported kernels,
$\CIc(M^2_\phi;\Hom(E)\otimes\Omega_R)$, where $\Omega_R$ is the fibre density
bundle lifted from the right factor (\cref{ch:product}); for $M$ compact this is
$\Psido{-\infty}(M/B;E)$.

\begin{remark}[Adjoining a trivial bundle]\label{scindex:rem:adjoin}
  We shall often need operators on $L^2(Z_b;E_b)\oplus\C^L$, for instance after
  stabilizing an elliptic family by finitely many sections as in
  \cref{ch:families}. These are covered by the formalism: $M\sqcup B\to B$ is a
  fibration whose fibres are $Z_b\sqcup\{b\}$, and $\C^L$ is the space of sections
  of the trivial bundle of rank $L$ over the point fibre $\{b\}$. Smoothing
  families on $M\sqcup B$ are block matrices whose off-diagonal entries are given
  by finitely many smooth sections of $E$ and $E^*\otimes\Omega$, and whose
  lower right entry is a smooth map $B\to M(L,\C)$. Pseudodifferential families
  on $(M\sqcup B)/B$ of order $m$ are matrices $\bigl(\begin{smallmatrix}P&S\\
  T&c\end{smallmatrix}\bigr)$ with $P\in\Psido m(M/B;E)$ and $S$, $T$, $c$ as
  just described; they compose as matrices and the smoothing families form an
  ideal. We use this without further comment.
\end{remark}

\subsection{K-theory}
We use the conventions of \cref{ch:bott,ch:families}. For a manifold $X$,
$\Kc(X)=\Kc^0(X)$ is the group of equivalence classes of triples
$(E_+,E_-,a)$, with $a:E_+\to E_-$ an isomorphism outside a compact set, under
bundle isomorphisms, change of $a$ on compact sets, homotopies and
stabilization. If $X$ is the total space of a vector bundle over $M$, bundles
over $X$ are pulled back from $M$ (\cref{bott:lem:bundle-homotopy}), so every class is
represented by a triple $(\varpi^*E_+,\varpi^*E_-,a)$ with $E_\pm$ bundles over
$M$ and $a\in\CI(X;\Hom(\varpi^*E_+,\varpi^*E_-))$ defined everywhere and
invertible outside a compact set; we usually drop $\varpi^*$. We also use:
\begin{enumerate}[label=(K\arabic*),leftmargin=*]
\item\label{scindex:K-Ginf} For a compact manifold $Y$ of positive dimension
  (not necessarily connected) and a Hermitian bundle $F\to Y$,
  $\Ginf(Y;F)\subset\Id+\Psido{-\infty}(Y;F)$ is the group of invertible
  perturbations of the identity by smoothing operators, and for compact $X$
  \begin{equation*}
    K^{-1}(X)=[X;\Ginf(Y;F)],\qquad K^{-2}(X)=[X\times\Sphere,X\times\{1\};
    \Ginf(Y;F),\Id],
  \end{equation*}
  smooth homotopy classes of smooth maps, respectively of pointed maps
  (\ref{bott:B:K1} and \cref{bott:def:K2}). By the finite rank approximation
  \ref{bott:B:finite} and \cref{bott:lem:stable}, every class is represented by a
  map into a subgroup $\GL(V)=\{\Id-\Pi_V+g\Pi_V\}$, $V\subset\CI(Y;F)$
  finite dimensional, and two such maps define the same class if and only if
  they are stably homotopic as maps into $\GL(N,\C)$ (after choosing orthonormal
  bases). In this sense the groups do not depend on $(Y,F)$.
\item\label{scindex:K-toeplitz} (Bott periodicity, first form;
  \cref{bott:thm:index-iso}.) The Toeplitz families index
  (\cref{bott:def:families-index}) is a natural isomorphism
  $\indT:K^{-2}(X)\to K^0(X)$ for compact $X$. If $p:X\to M(N,\C)$ is a smooth
  family of projections then, by \eqref{bott:eq:bott-map},
  \begin{equation}\label{scindex:eq:bott-map}
    \indT\bigl[pz^{-1}+(\Id-p)\bigr]=[\Ran p],\qquad
    \indT\bigl[pz+(\Id-p)\bigr]=-[\Ran p],
  \end{equation}
  where $z=e^{i\theta}$ is the variable on $\Sphere$.
\item\label{scindex:K-families} From \cref{ch:families}: the exterior product
  $\boxtimes$ and the product $x\cdot y$ for bundles over the same base, defined by
  the Clifford product of symbols (\cref{families:def:product}); the extension
  maps $j_*$ for open embeddings (\ref{families:K-open}); the Thom class
  $\lambda_{\mathcal V}\in\Kc(T^*(\mathcal V/M))$ of a real vector bundle
  \eqref{families:eq:thom}, represented on $\Lambda W$,
  $W=\varpi^*(\mathcal V\otimes\C)$, by Clifford multiplication
  $c(w)=w\wedge\cdot-\iota_w$, $w=v+i\eta$ (in the notation of
  \cref{bott:def:bott-class}, $\lambda_{\mathcal V}=\beta_{\mathcal V\otimes\C}$);
  the Thom homomorphism $\tau(x)=\lambda_{\mathcal V}\cdot\varpi^*x$; the Bott
  element $b_1\in\Kc(T^*I)$, the maps $\beta_N(y)=b_1\boxtimes\dots\boxtimes
  b_1\boxtimes y$, which are isomorphisms by \ref{families:BP}; the
  identification $\lambda_\R=b_1$ under $I\cong\R$
  (\cref{families:rem:bott-generator}) and hence $\lambda_{\R^N}=
  b_1^{\boxtimes N}$; and the topological index \eqref{families:eq:indt},
  $\indt=\beta_N^{-1}\circ j_*\circ\Theta_*\circ\tau$, defined by an embedding
  $M\hookrightarrow\R^N\times B$ with tubular neighbourhood
  $\Theta:\mathcal V\cong\Omega\subset I^N\times B$
  (\cref{families:lem:embedding}).
\item\label{scindex:K-inda} The analytic index $\inda:\Kc(T^*(M/B))\to K^0(B)$
  for a fibration of compact manifolds (\cref{product:def:index-bundle},
  \cref{families:lem:stabilization}): if $A\in\Psido0(M/B;E_+,E_-)$ is elliptic
  and $S:\C^L\to\CI(M;E_-)$ is given by $L$ sections such that $A\oplus S$ is
  surjective on every fibre, then the null spaces of $A\oplus S$ form a smooth
  bundle $\mathcal N(A,S)$ over $B$ and $\inda[\sigma_0(A)]=
  [\mathcal N(A,S)]-[\C^L]$.
\end{enumerate}

\subsection{The semiclassical calculus}
We use the semiclassical calculus of \cref{ch:semiclassical}, in the version for
families on the fibres of a fibration (\cref{semiclassical:rem:parameters}). For
$q\in\CI([0,1]_\epsilon;\Sym m(\R^n;\R^n))$ the semiclassical (left)
quantization is (\cref{semiclassical:def:psisl})
\begin{equation}\label{scindex:eq:Ope}
  \Ope(q)u(z)=(2\pi\epsilon)^{-n}\int e^{i(z-z')\cdot\zeta/\epsilon}
  q(\epsilon,z,\zeta)u(z')\,dz'd\zeta
  =\int e^{i(z-z')\cdot\xi}q(\epsilon,z,\epsilon\xi)u(z')\,dz'\,\dbar\xi,
\end{equation}
$\epsilon\in(0,1]$. Its kernel is $\epsilon^{-n}\check q(\epsilon,z,(z-z')/
\epsilon)$, where $\check q$ is the inverse Fourier transform in the last
variable (\eqref{semiclassical:eq:kernel}). If $q$ is Schwartz in $\zeta$ the
kernel concentrates, as $\epsilon\downarrow0$, in an $\epsilon$-neighbourhood of
the diagonal. The results of \cref{ch:semiclassical} are stated for compact
manifolds; since all constructions are local
(\cref{semiclassical:lem:structure}), they hold for families with kernels
supported over a fixed compact subset of a non-compact fibred manifold, which is
the form in which we state them.

\begin{enumerate}[label=(S\arabic*),leftmargin=*]
\item\label{scindex:S-smoothing} (Semiclassical smoothing families;
  \cref{semiclassical:prop:smoothing-X}.) For a fibration $\phi:M\to B$ and
  $E\to M$, $\Psisl{-\infty}(M/B;E)$ is the space of families $A=A(\epsilon)$,
  $\epsilon\in(0,1]$, of elements of $\Psido{-\infty}_{\mathrm c}(M/B;E)$ whose
  kernels are supported in $(0,1]\times K\times_BK$ for some compact
  $K\subset M$, and which are semiclassical: away from $\Diag$ the kernel lies in
  $\epsilon^\infty\CI([0,1]\times M^2_\phi)$, and in local coordinates $(y,z)$,
  $y$ in the base and $z$ in the fibre, near $\Diag$ it is
  $\epsilon^{-n}K(\epsilon,y,z,(z-z')/\epsilon)$ (times the coordinate density)
  with $K$ smooth and rapidly decreasing, with all derivatives, in the last
  variable, $n=\dim Z$. The families whose kernels are smooth on
  $[0,1]\times M^2_\phi$ and vanish to infinite order at $\epsilon=0$ form the
  ideal $\Psiinf(M/B;E)$, denoted $\epsilon^\infty\Psisl{-\infty}=
  \dCI([0,1];\Psido{-\infty})$ in \cref{ch:semiclassical}.
\item\label{scindex:S-symbol} (The semiclassical symbol;
  \cref{semiclassical:prop:smoothing-X,semiclassical:thm:manifold-calculus}.)
  There is a surjective linear map
  \begin{equation*}
    \sigsl:\Psisl{-\infty}(M/B;E)\longrightarrow\Sc(T^*(M/B);\End E),
  \end{equation*}
  locally $\sigsl(A)(y,z,\zeta)=\hat K(0,y,z,\zeta)$, onto the smooth sections
  which are Schwartz in the fibres of $T^*(M/B)$ and supported over a compact
  subset of $M$. It is multiplicative, $\sigsl(AB)=\sigsl(A)\sigsl(B)$, and
  $\sigsl(A^*)=\sigsl(A)^*$. Its null space is
  $\epsilon\Psisl{-\infty}(M/B;E)$, and
  $\bigcap_k\epsilon^k\Psisl{-\infty}=\Psiinf$. Every sequence
  $A_k\in\epsilon^k\Psisl{-\infty}$, with kernels supported in a fixed compact
  set, has an asymptotic sum $A\in\Psisl{-\infty}$,
  $A-\sum_{k<K}A_k\in\epsilon^K\Psisl{-\infty}$ for all $K$. A smooth bundle
  map $e\in\CI(M;\End E)$, regarded as a multiplication operator, satisfies
  $eA,Ae\in\Psisl{-\infty}$ with symbols $e\sigsl(A)$, $\sigsl(A)e$. If
  $\sigsl(A)$ is supported in $\varpi^{-1}(K)$, $K\subset M$ compact, then $A$
  may be chosen (modifying it by an element of $\Psiinf$) with kernel supported
  in any given neighbourhood of $\Diag\cap K\times_BK$
  (\cref{semiclassical:cor:localization}). All of this holds with smooth
  dependence on additional parameters (\cref{semiclassical:rem:parameters}).
\item\label{scindex:S-invariance} (Naturality;
  \cref{semiclassical:prop:invariance}.) A diffeomorphism $F:M\to M'$ of
  fibrations over $B$, covered by a unitary bundle isomorphism $E\to E'$,
  conjugates $\Psisl{-\infty}(M/B;E)$ to $\Psisl{-\infty}(M'/B;E')$ (with the
  obvious density factors), and the symbol transforms by the induced bundle
  isomorphism $T^*(M/B)\to T^*(M'/B)$.
\item\label{scindex:S-order0} (Finite order; compact fibres;
  \cref{semiclassical:thm:manifold-calculus}.) For a fibration of compact
  manifolds, $\Psisl{m}(M/B;E_+,E_-)$ is the space of families which are,
  locally, $\Ope(q)$ with $q$ in $\Ssl m=\CI([0,1];\Sym m)$, plus elements
  of $\Psiinf$. Then $\Psisl{-\infty}=\bigcap_k\Psisl{-k}$, the space
  $\Psisl0$ is an algebra
  containing $\Psisl{-\infty}$ as an ideal, and for fixed $\epsilon>0$,
  $A(\epsilon)\in\Psido0(M/B)$. There are two multiplicative symbols: the
  semiclassical symbol $\sigsl(A)$, a symbol of order $0$ on $T^*(M/B)$ with
  values in $\Hom(E_+,E_-)$, locally $q(0,\cdot)$, and the uniform principal symbol $\tilde\sigma_0(A)$, the class of
  $q$ modulo $\Ssl{-1}$; for $\epsilon>0$ the principal symbol of $A(\epsilon)$ is
  represented by $q(\epsilon,z,\epsilon\xi)$, by part (f) of
  \cref{semiclassical:thm:manifold-calculus}. For classical symbols which are
  homogeneous of degree $0$ for $\abs\zeta\ge1$ the rescaling does not change the
  homogeneous part, so $\sigma_0(A(\epsilon))$ is the boundary value of
  $q(\epsilon,\cdot)$ at infinity. The sequences
  \begin{equation*}
    \epsilon\Psisl{m}\hookrightarrow\Psisl m\xrightarrow{\sigsl}\Sym m,\qquad
    \Psisl{m-1}\hookrightarrow\Psisl m\xrightarrow{\tilde\sigma_m}\Ssl m/\Ssl{m-1}
  \end{equation*}
  are exact (\cref{semiclassical:prop:symbol-maps}); if $T_k\in
  \epsilon^k\Psisl{-k}$ then there is $T\in\Psisl0$ with
  $T-\sum_{k<K}T_k\in\epsilon^K\Psisl{-K}$ for all $K$, and
  $\bigcap_K\epsilon^K\Psisl{-K}=\Psiinf$; and every symbol $a$ of order $0$ on $T^*(M/B)$ with values in
  $\Hom(E_+,E_-)$ is $\sigsl(A)$ for some $A$, which may be taken with
  $\tilde\sigma_0(A)=[a]$ independent of $\epsilon$ (the quantization map $q_\epsilon$
  of \cref{semiclassical:thm:manifold-calculus}(d) applied to $a$, regarded as
  constant in $\epsilon$).
\item\label{scindex:S-euclid} (Euclidean estimates;
  \cref{semiclassical:thm:composition,semiclassical:thm:L2}.) On $\R^k$, for
  symbols $p,q\in\Sym0(\R^k;\R^k)$ (bounded with all derivatives in the base
  variable), $\sup_{0<\epsilon\le1}\norm{\Ope(p)}_{L^2}\le C$, and
  \begin{equation*}
    \norm{\Ope(p)\Ope(q)-\Ope(pq)}_{L^2}\le C\epsilon,\quad 0<\epsilon\le1,
  \end{equation*}
  since $p\#_\epsilon q-pq\in\epsilon\Ssl{-1}$. Moreover
  $\Ope(p)\Ope(q)=\Ope(p\#_\epsilon q)$ with
  \begin{equation*}
    p\#_\epsilon q\sim\sum_\alpha\frac{\epsilon^{\abs\alpha}}{\alpha!}
    \pa_\zeta^\alpha p\,D_z^\alpha q
  \end{equation*}
  by \eqref{semiclassical:eq:composition}; the expansion is \emph{exact} (a finite
  sum) if $q$ is a polynomial in $z$, or if $p$ is a polynomial in $\zeta$, and in
  these cases it holds for polynomial symbols of any order.
\end{enumerate}

We add two consequences. For compact fibres they are
\cref{semiclassical:thm:L2-X} and \cref{semiclassical:cor:G}; we include the short
proofs, which apply verbatim to compactly supported families on non-compact fibres.

\begin{lemma}\label{scindex:lem:uniform}
  If $A\in\Psisl{-\infty}(M/B;E)$ then $\sup_{\epsilon,b}\norm{A_b(\epsilon)}
  _{L^2}<\infty$, and if $A\in\Psiinf(M/B;E)$ then $\sup_b\norm{A_b(\epsilon)}
  =O(\epsilon^\infty)$.
\end{lemma}
\begin{proof}
  The second statement is clear, since the kernels are smooth with compact
  support and all their sup norms are $O(\epsilon^\infty)$. For the first,
  by a partition of unity we may assume the kernel is
  $\epsilon^{-n}K(\epsilon,y,z,(z-z')/\epsilon)$ with compact support in $z$.
  Since $\sup_{\epsilon,y,z}\abs{K(\epsilon,y,z,Z)}\le C\ang Z^{-n-1}$,
  \begin{equation*}
    \sup_z\int\epsilon^{-n}\abs{K(\epsilon,y,z,(z-z')/\epsilon)}dz'
    \le C\int\ang Z^{-n-1}dZ,
  \end{equation*}
  and the same holds with the roles of $z$ and $z'$ exchanged. Schur's test
  gives the bound.
\end{proof}

\begin{lemma}[Invertibility]\label{scindex:lem:invertible}
  Let $A\in\Psisl{-\infty}(M/B;E)$, possibly depending smoothly on parameters
  in a compact manifold $P$, and suppose that $\Id+\sigsl(A)$ is invertible at
  every point. Then there is $\epsilon_0>0$ such that $\Id+A_b(\epsilon)$ is
  invertible on $L^2(Z_b;E_b)$ for all $b\in B$, all parameters and all
  $\epsilon\in(0,\epsilon_0]$, and $(\Id+A)^{-1}=\Id+A'$ with
  $A'\in\Psisl{-\infty}(M/B;E)$ (restricted to $\epsilon\le\epsilon_0$),
  depending smoothly on the parameters. If $P$ is a complex manifold and $A$
  depends holomorphically on the parameters, so does $A'$.
\end{lemma}
\begin{proof}
  The symbol $a'=(\Id+\sigsl(A))^{-1}-\Id$ is again Schwartz and supported over a
  compact set. Choose $A'_0$ with $\sigsl(A'_0)=a'$. Then
  $(\Id+A)(\Id+A'_0)=\Id+\epsilon R_1$, $R_1\in\Psisl{-\infty}$. Inductively, if
  $(\Id+A)(\Id+A'_{(k)})=\Id+\epsilon^kR_k$, set
  $A'_{(k+1)}=A'_{(k)}-\epsilon^kQ_k$ with $\sigsl(Q_k)=(\Id+\sigsl(A))^{-1}
  \sigsl(R_k)$; then the error is $O(\epsilon^{k+1})$. Taking an asymptotic sum
  gives $A''$ with $(\Id+A)(\Id+A'')=\Id+R$, $R\in\Psiinf$, and similarly a
  left inverse modulo $\Psiinf$. By \cref{scindex:lem:uniform},
  $\norm R=O(\epsilon^\infty)$, so for small $\epsilon$, $\Id+R$ is invertible
  and
  \begin{equation*}
    (\Id+R)^{-1}=\Id-R+R(\Id+R)^{-1}R.
  \end{equation*}
  The last term has kernel $(w,w')\mapsto\inner{(\Id+R)^{-1}R(\cdot,w')}
  {R^*(w,\cdot)}$, which is smooth, with all derivatives (in $w$, $w'$,
  $\epsilon$ and the parameters) $O(\epsilon^\infty)$, because derivatives of
  $(\Id+R)^{-1}$ are products of $(\Id+R)^{-1}$ with derivatives of $R$. So
  $(\Id+R)^{-1}\in\Id+\Psiinf$, $\Id+A$ has a right inverse in
  $\Id+\Psisl{-\infty}$, and in the same way a left inverse; they coincide.
  Holomorphy is preserved at each step (the asymptotic sum may be taken with
  coefficients holomorphic in the parameters, and inversion is holomorphic).
\end{proof}

\begin{remark}[The semiclassical cotangent bundle]\label{scindex:rem:slT}
  The symbol $\sigsl(A)$ depends on the parameter $\epsilon$, not only on the
  family of operators: under a reparametrization $\epsilon'=\epsilon T(\epsilon)$
  with $T>0$ smooth, the symbol is composed with the fibre dilation by $T(0)$.
  Invariantly, $\sigsl(A)$ is a function on a bundle ${}^{\mathrm{sl}}T^*(M/B)$
  which is isomorphic, but not canonically, to $T^*(M/B)$; a covector of
  ${}^{\mathrm{sl}}T^*$ at $z$ is the class of $df/\epsilon$ for a function $f$
  with $f(z)=0$. Once $\epsilon$ is fixed, as it is throughout this chapter, the
  identification is fixed. Since the dilations are homotopic to the identity,
  none of the K-theoretic statements below depends on this choice; see
  \cref{semiclassical:rem:sl-cotangent}.
\end{remark}

\section{Smoothing families and the K-theory of the base}\label{scindex:sec:smoothing}

Let $\phi:M\to B$ be a fibration over the compact manifold $B$ and $E\to M$ a
Hermitian bundle. Set
\begin{equation*}
  \begin{aligned}
  \Ginf_{\mathrm c}(M/B;E)=\bigl\{g=\Id+A;\ &A\in\Psido{-\infty}_{\mathrm c}(M/B;E),\\
  &\Id+A_b\text{ invertible on }L^2(Z_b;E_b)\text{ for every }b\in B\bigr\}.
  \end{aligned}
\end{equation*}
If $g=\Id+A\in\Ginf_{\mathrm c}$, then $g_b^{-1}=\Id+A'_b$ with $A'$ again in
$\Psido{-\infty}_{\mathrm c}(M/B;E)$: indeed $A'=-A+A(\Id+A)^{-1}A$, and as in the
proof of \cref{scindex:lem:invertible} the last term has a smooth kernel
depending smoothly on $b$, with support in that of $A$. Thus $\Ginf_{\mathrm c}
(M/B;E)$ is a group, the group of smooth sections of a bundle of groups over $B$
with fibre $\Ginf_{\mathrm c}(Z_b;E_b)$. We assign to each of its elements a class
in $K^{-1}(B)$. (Melrose does this by showing that the Hilbert bundle with fibres
$L^2(Z_b;E_b)$ has a smooth orthonormal frame; see \cref{scindex:ex:frame}. The
following embedding is more elementary.)

\subsection{Admissible embeddings}
\begin{definition}\label{scindex:def:admissible}
  Let $K\subset M$ be compact and $W\supset K$ open. An \emph{admissible
  embedding} of $(M/B,E)$ over $W$ is a family of linear isometries
  \begin{equation*}
    J_b:L^2(W_b;E_b)\longrightarrow L^2(Y;F),\quad W_b=W\cap Z_b,
  \end{equation*}
  where $Y$ is a compact manifold and $F\to Y$ a Hermitian bundle, such that
  $J_b$ maps $\CIc(W_b;E_b)$ into $\CI(Y;F)$ and, for every
  $A\in\Psido{-\infty}_{\mathrm c}(M/B;E)$ with kernel supported in
  $K\times_BK$, the map $b\mapsto J_bA_bJ_b^*$ is a smooth map
  $B\to\Psido{-\infty}(Y;F)$ (depending continuously on $A$).
\end{definition}

\begin{lemma}\label{scindex:lem:embedding}
  For every compact $K\subset M$ there is an admissible embedding of $(M/B,E)$
  over some $W\supset K$, with $Y=\Torus^N$ and $F=\C^R$ for some $N$, $R$.
\end{lemma}
\begin{proof}
  Let $\bar M\supset M$ be the compact fibration of which $M$ is an open
  subset. As in \cref{families:lem:embedding}, choose an embedding $\iota:\bar M\to\R^N$ such
  that $(\iota,\bar\phi)$ embeds $\bar M$ into $\R^N\times B$, with image in
  $(0,2\pi)^N\times B$, and let $\mathcal V_m=d\iota(T_m(\bar M/B))^\perp$. For
  some $\delta>0$ the map $\Theta(m,v)=\iota(m)+v$ is, for each $b$, a
  diffeomorphism from $\{(m,v);\ m\in\bar Z_b,\ v\in\mathcal V_m,\ \abs v<\delta\}$
  onto an open set $\Omega_b\subset(0,2\pi)^N\subset\Torus^N$, depending smoothly
  on $b$; let $p_b:\Omega_b\to\bar Z_b$ be the corresponding projection. Choose
  $h\in\CIc([0,\delta))$ real, not identically zero near $0$, and define
  $\gamma'(y)=h(\abs v)$ for $y=\Theta(m,v)$, extended by $0$. By the
  coarea formula
  \begin{equation*}
    \int_{\Omega_b}\abs{\gamma'(y)}^2\abs{u(p_b(y))}^2dy=\int_{\bar Z_b}
    f(m)\abs{u(m)}^2\nu(m)
  \end{equation*}
  for a smooth positive function $f$ on $\bar M$ (the integral of
  $\abs{\gamma'}^2$ over $p_b^{-1}(m)$, with respect to the measure induced by
  $dy$ and $\nu$). Set $\gamma=\gamma'\cdot(f\circ p)^{-1/2}$. Next let
  $W\supset K$ be open with compact closure in $M$, and $W'\supset\bar W$ open
  with compact closure; as in the proof of \cref{bott:lem:bundles},
  a finite cover of $\bar W'$ by trivializing open sets and a partition of
  unity give a fibrewise injective bundle map $E|_{W'}\to\C^R$, which after
  Gram--Schmidt we may take isometric, $\kappa:E|_{W'}\to\C^R$. Define
  \begin{equation*}
    (J_bu)(y)=\gamma(y)\,\kappa\bigl(u(p_b(y))\bigr)\ \text{ if }p_b(y)\in W,
    \qquad(J_bu)(y)=0\ \text{ otherwise}.
  \end{equation*}
  By construction $\norm{J_bu}=\norm u$. If $u\in\CIc(W_b)$ then $J_bu$ is
  smooth with compact support in $\Omega_b$. For $A$ with kernel $A(m,m')$
  supported in $K\times_BK$, $J_bA_bJ_b^*$ has the kernel
  \begin{equation*}
    \gamma(y)\,\kappa\,A_b(p_b(y),p_b(y'))\,\kappa^*\,\gamma(y')\,j_b(y'),
  \end{equation*}
  where $j_b$ is the smooth positive Jacobian factor relating $\nu$, $dy$ and the
  measures on the fibres of $p_b$; it is smooth in $(b,y,y')$ with
  compact support in $\Omega_b\times\Omega_b$.
\end{proof}

For a non-compact vector bundle fibre we shall also need kernels which are
Schwartz, rather than compactly supported, along the fibres.

\begin{remark}\label{scindex:rem:schwartz}
  Let $\varpi:\mathcal V\to M$ be a real vector bundle with a metric over a
  fibration of compact manifolds $M\to B$. Let $\Psido{-\infty}_{\mathcal S}
  (\mathcal V/B;E)$ denote the families whose kernels are smooth on
  $\mathcal V\times_B\mathcal V$ and Schwartz in the fibre variables, uniformly
  with all derivatives (i.e.\ they extend smoothly, vanishing to infinite order at
  the boundary, to the fibrewise radial compactification; see
  \cref{ch:distributions}). Composing the construction above with the fibre
  compression $v\mapsto v(1+\abs v^2)^{-1/2}$, which identifies $\mathcal V$ with
  the open unit ball bundle, and using a function $\gamma$ which is Gaussian in
  the original variable, one obtains isometries $J_b:L^2(\mathcal V_b)\to
  L^2(\Torus^N;\C^R)$ with the properties of \cref{scindex:def:admissible} for
  all elements of $\Psido{-\infty}_{\mathcal S}$; the Jacobian factors grow at
  most polynomially in $\abs v$ and are absorbed by the rapid decay. Everything
  below applies to such families without change.
\end{remark}

\subsection{The class of a section}
\begin{definition}\label{scindex:def:class}
  Let $g=\Id+A\in\Ginf_{\mathrm c}(M/B;E)$, with $A$ supported in $K\times_BK$,
  and let $J$ be an admissible embedding over $W\supset K$. Then
  \begin{equation*}
    \rho_J(g)(b)=\Id+J_bA_bJ_b^*\in\Ginf(Y;F),
  \end{equation*}
  and we set $[g]=[\rho_J(g)]\in[B;\Ginf(Y;F)]=K^{-1}(B)$.
\end{definition}
Indeed $\rho_J(g)(b)$ is invertible: with respect to $L^2(Y;F)=\Ran J_b\oplus
(\Ran J_b)^\perp$ it is $J_bg_bJ_b^*\oplus\Id$, where $g_b$ acts on $L^2(W_b)$ (it
preserves this space and its complement, being the identity on
$L^2(Z_b\setminus K)$), and its inverse is $\Id+J_bA'_bJ_b^*$. Note that
$J^*J=\Id$ implies $\rho_J(gh)=\rho_J(g)\rho_J(h)$.

\begin{lemma}\label{scindex:lem:class}
  \begin{enumerate}[label=(\roman*)]
  \item If $I:L^2(Y;F)\to L^2(Y';F')$ is an isometry which maps $\CI(Y;F)$ into
    $\CI(Y';F')$ and $\Psido{-\infty}(Y;F)$ continuously into
    $\Psido{-\infty}(Y';F')$ by $T\mapsto ITI^*$, then the map
    $h\mapsto IhI^*+\Id-II^*$ induces the canonical identification
    $[B;\Ginf(Y;F)]\cong[B;\Ginf(Y';F')]$ of \ref{scindex:K-Ginf}.
  \item The class $[g]\in K^{-1}(B)$ does not depend on $K$, $W$ or $J$.
  \item $[gh]=[g]+[h]$; smoothly homotopic sections define the same class; and if
    $U\subset M$ is open and $g\in\Ginf_{\mathrm c}(U/B;E)$, then the class of $g$
    is the same as that of its extension (by the identity) to an element of
    $\Ginf_{\mathrm c}(M/B;E)$.
  \item If $\Phi:E\to E'$ is a unitary isomorphism covering a diffeomorphism of
    fibrations over $B$ (in particular $\Phi$ may be a unitary bundle map on
    $M$), then $[\Phi g\Phi^{-1}]=[g]$.
  \end{enumerate}
\end{lemma}
\begin{proof}
  (i) By \ref{scindex:K-Ginf} a class is represented by $h$ with values in a
  subgroup $\GL(V)$, $V\subset\CI(Y;F)$ finite dimensional. Then $IhI^*+\Id-II^*$
  takes values in $\GL(IV)$, $IV\subset\CI(Y';F')$, and in the orthonormal bases
  $\{e_j\}$ of $V$ and $\{Ie_j\}$ of $IV$ the two maps are given by the same
  matrices. This is the identification of \ref{scindex:K-Ginf}.

  (ii) Enlarging $K$ does not change $\rho_J(g)$. If $J$, $J'$ are admissible over
  $W$, $W'$, both restrict to admissible embeddings over $W\cap W'\supset K$,
  and for $A$ supported in $K\times_BK$ the restriction does not change
  $JAJ^*$. So we may assume $W=W'$. Consider $Y\sqcup Y'$ with the bundle
  $F\sqcup F'$, so $L^2(Y\sqcup Y')=L^2(Y)\oplus L^2(Y')$, and
  \begin{equation*}
    R_{t,b}=(\cos t\,J_b,\ \sin t\,J'_b):L^2(W_b;E_b)\longrightarrow L^2(Y;F)\oplus
    L^2(Y';F'),\quad t\in[0,\tfrac\pi2].
  \end{equation*}
  Since the ranges of $(J_b,0)$ and $(0,J'_b)$ are orthogonal, $R_t^*R_t=
  \cos^2t\,J^*J+\sin^2t\,J'^*J'=\Id$, so each $R_t$ is admissible, and
  $\rho_{R_t}(g)$ is a homotopy in $\Ginf(Y\sqcup Y')$ from
  $\rho_J(g)\oplus\Id$ to $\Id\oplus\rho_{J'}(g)$. By (i), applied to the
  inclusions of $L^2(Y)$ and $L^2(Y')$ as summands, these represent $[\rho_J(g)]$
  and $[\rho_{J'}(g)]$.

  (iii) The first statement follows from $\rho_J(gh)=\rho_J(g)\rho_J(h)$. A
  homotopy is a section over $B\times[0,1]$, and one admissible embedding
  (\cref{scindex:lem:embedding} over $B\times[0,1]$) serves for all $t$. An
  admissible embedding of $M$ over $W\subset U$ is one of $U$.

  (iv) If $J$ is an admissible embedding for $(M'/B,E')$ over $W'$, then $J\Phi$
  is one for $(M/B,E)$ over the preimage of $W'$, and since $\Phi$ is unitary,
  $J(\Phi A\Phi^{-1})J^*=(J\Phi)A(J\Phi)^*$. Thus $\rho_J(\Phi g\Phi^{-1})=
  \rho_{J\Phi}(g)$, and (ii) applies.
\end{proof}

Applying this to the fibration $M\times\Sphere\to B\times\Sphere$, a smooth
\emph{pointed loop} $g:\Sphere\to\Ginf_{\mathrm c}(M/B;E)$, i.e.\ an element of
$\Ginf_{\mathrm c}((M\times\Sphere)/(B\times\Sphere);E)$ equal to $\Id$ over
$B\times\{1\}$, has a class $[g]\in K^{-2}(B)$ (take $J$ independent of $z$).

\subsection{The relative index of idempotents}
Let $p_0\in\CI(M;\End E)$ be a smooth family of idempotents ($p_0^2=p_0$), acting
as a multiplication operator.

\begin{definition}\label{scindex:def:d}
  A \emph{compactly supported idempotent perturbation} of $p_0$ is a smooth family
  $e=p_0+Q$, $Q\in\Psido{-\infty}_{\mathrm c}(M/B;E)$ (or, in the situation of
  \cref{scindex:rem:schwartz}, $Q\in\Psido{-\infty}_{\mathcal S}$), with
  $e_b^2=e_b$ for all $b$. Its \emph{relative index} is
  \begin{equation*}
    d(e,p_0)=\indT[g_{e,p_0}]\in K^0(B),\qquad
    g_{e,p_0}(z)=\bigl(ez^{-1}+\Id-e\bigr)\bigl(p_0z+\Id-p_0\bigr),\ z\in\Sphere.
  \end{equation*}
\end{definition}

For an idempotent $f$, $z\mapsto fz^{-1}+\Id-f$ is a loop of invertible operators
with inverse $fz+\Id-f$. Expanding, with $e=p_0+Q$, and using $p_0(\Id-p_0)=0$,
\begin{equation}\label{scindex:eq:loop-expansion}
  g_{e,p_0}(z)=\Id+Q(2p_0-\Id)+z^{-1}Q(\Id-p_0)-zQp_0,
\end{equation}
so $g_{e,p_0}$ is a pointed loop in $\Ginf_{\mathrm c}(M/B;E)$, and $d(e,p_0)$ is
defined. (The construction is that of the Bott map \eqref{scindex:eq:bott-map},
applied to $e$, followed by subtraction of the same construction for $p_0$; the
two infinite rank contributions cancel.)

\begin{proposition}\label{scindex:prop:d}
  \begin{enumerate}[label=(d\arabic*)]
  \item\label{scindex:d-homotopy} If $e_t$, $t\in[0,1]$, is a smooth family of
    compactly supported idempotent perturbations of $p_0$ (supports in a fixed
    compact set), then $d(e_0,p_0)=d(e_1,p_0)$.
  \item\label{scindex:d-sum} $d(e\oplus e',p_0\oplus p_0')=d(e,p_0)+d(e',p'_0)$,
    and $d(p_0,p_0)=0$.
  \item\label{scindex:d-conj} If $\Phi$ is as in \cref{scindex:lem:class}(iv) and
    commutes with $p_0$ (resp.\ carries $p_0$ to $p_0'$), then
    $d(\Phi e\Phi^{-1},p'_0)=d(e,p_0)$. This also holds for $\Phi$ a smooth family of
    invertible (not necessarily unitary) bundle maps on $M$.
  \item\label{scindex:d-finite} If $\Pi\in\Psido{-\infty}_{\mathrm c}(M/B;E)$ is
    a smooth family of finite rank idempotents with $p_0\Pi=\Pi p_0=0$, then the
    ranges of $\Pi_b$ form a smooth vector bundle $\Ran\Pi$ over $B$ and
    $d(p_0+\Pi,p_0)=[\Ran\Pi]$. If instead $p_0\Pi=\Pi p_0=\Pi$ then
    $d(p_0-\Pi,p_0)=-[\Ran\Pi]$.
  \item\label{scindex:d-module} If $F\to B$ is a vector bundle, lifted to $M$,
    then $d(e\otimes\Id_F,p_0\otimes\Id_F)=d(e,p_0)\cdot[F]$.
  \item\label{scindex:d-excision} If $U\subset M$ is open and $e-p_0$ is supported
    in $K\times_BK$, $K\subset U$ compact, then $d(e,p_0)$ is the same whether
    computed on $U$ or on $M$.
  \end{enumerate}
\end{proposition}
\begin{proof}
  \ref{scindex:d-homotopy}, \ref{scindex:d-sum} and \ref{scindex:d-excision}
  follow from \cref{scindex:lem:class}(iii), using that $\indT$ is a homomorphism,
  that block sums of loops represent sums in $K^{-2}(B)$, and that
  $g_{p_0,p_0}=\Id$. For \ref{scindex:d-conj}, $g_{\Phi e\Phi^{-1},p'_0}=\Phi
  g_{e,p_0}\Phi^{-1}$; if $\Phi$ is not unitary, the polar decomposition
  $\Phi_s=\Phi(\Phi^*\Phi)^{-s/2}$, $s\in[0,1]$, joins it to a unitary through
  invertible bundle maps, each of which commutes with $p_0$ if $\Phi$ and
  $\Phi^*$ do; in general conjugation by $\Phi_s$ is a homotopy of loops in
  $\Ginf_{\mathrm c}$, and we use \cref{scindex:lem:class}(iii),(iv).

  \ref{scindex:d-finite}. The rank of $\Pi_b$ is $\Tr\Pi_b$, which is locally
  constant, and if $v_1,\dots,v_r$ span $\Ran\Pi_{b_0}$ then
  $\Pi_bv_1,\dots,\Pi_bv_r$ are smooth local frames near $b_0$; so $\Ran\Pi$ is a
  smooth bundle, of smooth sections. If $p_0\Pi=\Pi p_0=0$ then $p_0$ and $\Pi$
  are commuting idempotents with product zero, so
  $(p_0+\Pi)z^{-1}+\Id-p_0-\Pi=(p_0z^{-1}+\Id-p_0)(\Pi z^{-1}+\Id-\Pi)$, the
  factors commuting, and therefore $g_{p_0+\Pi,p_0}=\Pi z^{-1}+\Id-\Pi$. It
  remains to show that the class of this loop is the Bott class of the bundle
  $N=\Ran\Pi$, so that \eqref{scindex:eq:bott-map} applies. Transport by an
  admissible embedding $J$; then $J\Pi J^*$ is a smooth family of finite rank
  idempotents in $\Psido{-\infty}(Y;F)$ with range $JN\cong N$, and we write
  $\Pi$ for it. Let $P_N$ be the orthogonal projection onto $N$ (given locally by
  orthonormal frames, so a smooth family of finite rank smoothing operators). Since
  $\Pi P_N=P_N$ and $P_N\Pi=\Pi$, $\Pi_s=P_N+(1-s)P_N\Pi(\Id-P_N)$ is an
  idempotent with range $N$ for $s\in[0,1]$; this deforms $\Pi$ to $P_N$. By the
  finite rank approximation \ref{bott:B:finite} there is a finite dimensional
  $V\subset\CI(Y;F)$ with orthogonal projection $P_V$ such that
  $\norm{P_N-P_VP_NP_V}\le\frac12$, hence $\norm{(\Id-P_V)P_N}\le\frac12$, for
  all $b$. Then $T_s=\Id-s(\Id-P_V)$ is injective on $N_b$, so
  $N_s=T_sN$ is a smooth bundle isomorphic to $N$, and the orthogonal
  projections $P_{N_s}$ deform $P_N$ to the projection onto $N_1=P_VN\subset V$.
  The loop $P_{N_1}z^{-1}+\Id-P_{N_1}$ lies in $\GL(V)$, where it is the loop
  $pz^{-1}+\Id-p$ of \eqref{scindex:eq:bott-map} with $\Ran p=N_1\cong N$. Hence
  $d(p_0+\Pi,p_0)=[N]$. If $p_0\Pi=\Pi p_0=\Pi$, then similarly
  $(p_0-\Pi)z^{-1}+\Id-p_0+\Pi=(p_0z^{-1}+\Id-p_0)(\Pi z+\Id-\Pi)$, so
  $g_{p_0-\Pi,p_0}=\Pi z+\Id-\Pi$, and the second formula in
  \eqref{scindex:eq:bott-map} gives $-[\Ran\Pi]$.

  \ref{scindex:d-module}. Write $F=\Ran q$ for a family of projections
  $q:B\to M(r,\C)$. By \ref{scindex:d-sum} and \ref{scindex:d-conj},
  $d(e\otimes\Id_F,p_0\otimes\Id_F)$ is the Toeplitz index of the loop
  $h=g\otimes q+\Id\otimes(\Id-q)$ on $L^2\otimes\C^r$, $g=g_{e,p_0}$. By
  \ref{scindex:K-Ginf} we may assume that $g$ takes values in $\GL(n,\C)$, and
  then $h$ takes values in $\GL(nr,\C)$. Let $T_g$ be the Toeplitz family of $g$,
  $k$ admissible for it with complement $W$ (in the sense of \cref{bott:def:index-data}, \cref{bott:lem:index-bundle}).
  Since $q$ is constant on the circle it commutes with the Szeg\H{o} projector
  and with $P_k$, so $T_h=T_g\otimes q+\Id\otimes(\Id-q)$, $k$ is admissible for
  $T_h$, and $W\otimes F\oplus\Ran P_k\otimes\Ran(\Id-q)$ is a complement. Hence
  $\indT(h)=[\Ran P_k\otimes\C^r]-[W\otimes F]-[\Ran P_k\otimes\Ran(\Id-q)]
  =([\Ran P_k]-[W])[F]=\indT(g)[F]$.
\end{proof}

\begin{remark}\label{scindex:rem:relative-projections}
  If $e=P$, $p_0=Q$ are orthogonal projections, then $QP:\Ran P\to\Ran Q$ is a
  Fredholm family (it differs from the identity of $\Ran Q$ by a compact operator
  after identifying $\Ran P$ with $\Ran Q$), and $d(P,Q)$ is its index bundle, the
  \emph{relative index} of the pair. This is Melrose's description of the even
  semiclassical index (``the formal difference of the pair $P\ominus\pi_\infty$'');
  we do not need it, but see \cref{scindex:ex:relative}.
\end{remark}

\section{The semiclassical index}\label{scindex:sec:sl-index}

\subsection{An idempotent attached to an invertible-at-infinity map}
The following purely algebraic construction is used twice: for symbols, to
represent a compactly supported K-class by idempotents which are trivial outside
a compact set, and for operators, to express the index bundle of an elliptic
family as a relative index.

\begin{lemma}\label{scindex:lem:milnor}
  Let $A:H_+\to H_-$ and $B:H_-\to H_+$ be linear maps (between vector spaces,
  modules, or elements of a matrix algebra), and set $R_+=\Id-BA$,
  $R_-=\Id-AB$. Then
  \begin{equation*}
    W=\begin{pmatrix}A+R_-A&-R_-\\R_+&B\end{pmatrix}:H_+\oplus H_-\to H_-\oplus
    H_+,
  \end{equation*}
  \begin{equation*}
    W'=\begin{pmatrix}B&R_+\\-R_-&A+R_-A\end{pmatrix}:H_-\oplus H_+\to H_+\oplus
    H_-
  \end{equation*}
  are inverse to each other, and
  \begin{equation}\label{scindex:eq:milnor}
    e_{A,B}=W\begin{pmatrix}\Id&0\\0&0\end{pmatrix}W'=
    \begin{pmatrix}\Id-R_-^2&(\Id+R_-)AR_+\\R_+B&R_+^2\end{pmatrix}
  \end{equation}
  is an idempotent on $H_-\oplus H_+$. With $p_0=\bigl(\begin{smallmatrix}\Id&0\\0&0
  \end{smallmatrix}\bigr)$ (the projection onto $H_-$), $e_{A,B}-p_0$ is a
  matrix of polynomials in $A$, $B$, $R_\pm$ each term of which contains a factor
  $R_+$ or $R_-$. If $B=A^{-1}$ then $e_{A,B}=p_0$.
\end{lemma}
\begin{proof}
  Direct multiplication, using $AR_+=R_-A$ and $R_+B=BR_-$. For instance the
  upper left entry of $W'W$ is
  \begin{equation*}
    B(A+R_-A)+R_+^2=BA+R_+BA+R_+^2=BA+R_+=\Id,
  \end{equation*}
  and the lower left entry is
  \begin{equation*}
    -R_-(A+R_-A)+(A+R_-A)R_+=-R_-^2A+R_-AR_+=0;
  \end{equation*}
  the other
  entries, and those of $WW'$, are checked in the same way. Then
  \begin{equation*}
    e_{A,B}^2=W\diag(\Id,0)W'W\diag(\Id,0)W'=e_{A,B},
  \end{equation*}
  and \eqref{scindex:eq:milnor} follows
  from $(A+R_-A)B=(\Id+R_-)(\Id-R_-)$.
\end{proof}

The idempotent $e_{A,B}$ is the image of the class of $A$ under the boundary map
in algebraic K-theory, in the explicit form due to Milnor. For an invertible
$A$ it is trivial; in general it measures the failure of $B$ to be an inverse.

\subsection{Quantization of idempotents}
Let $\phi:M\to B$ be a fibration, $E\to M$, and $p_0\in\CI(M;\End E)$ an
idempotent.

\begin{proposition}\label{scindex:prop:quantize}
  Let $e\in\CI(T^*(M/B);\End E)$ be a family of idempotents such that $e-p_0\in
  \Sc(T^*(M/B);\End E)$. Then there are $\epsilon_0>0$ and $\mathfrak E=p_0+Q$,
  $Q\in\Psisl{-\infty}(M/B;E)$, with $\sigsl(Q)=e-p_0$ and $\mathfrak E(\epsilon)^2
  =\mathfrak E(\epsilon)$ for $0<\epsilon\le\epsilon_0$; we call $\mathfrak E$ a
  \emph{semiclassical quantization} of $e$. If $\mathfrak E'$ is another one, then
  for $\epsilon$ small $\mathfrak E(\epsilon)$ and $\mathfrak E'(\epsilon)$ are
  joined by a smooth family of compactly supported idempotent perturbations of
  $p_0$. The same holds with smooth dependence on parameters; in particular a
  smooth homotopy $e_t$, with $e_t-p_0$ supported in a fixed compact set, has
  quantizations $\mathfrak E_t$ depending smoothly on $t$.
\end{proposition}
\begin{proof}
  Choose $Q_0\in\Psisl{-\infty}$ with $\sigsl(Q_0)=e-p_0$ and set $P=p_0+Q_0$. Then
  $P^2-P=p_0Q_0+Q_0p_0+Q_0^2-Q_0\in\Psisl{-\infty}$ has symbol
  $(p_0+q)^2-(p_0+q)=e^2-e=0$, $q=e-p_0$, so $P^2-P=\epsilon E$ with
  $E\in\Psisl{-\infty}$. For $s$ on the circle $\Gamma=\{\abs{s-1}=\frac12\}$,
  $\abs{s(s-1)}\ge\frac14$ and
  \begin{equation*}
    (s-P)(s-1+P)=s(s-1)-(P^2-P)=s(s-1)\Bigl(\Id-\frac{\epsilon E}{s(s-1)}\Bigr).
  \end{equation*}
  The last factor is $\Id$ plus an element of $\Psisl{-\infty}$, depending
  holomorphically on $s$ near $\Gamma$, with symbol $0$; by
  \cref{scindex:lem:invertible} it is invertible for small $\epsilon$, with
  inverse $\Id+\epsilon E'_s$, $E'_s\in\Psisl{-\infty}$ holomorphic in $s$. Since
  $s-P$ and $s-1+P$ commute, $s\notin\Spec(P)$ and
  \begin{equation*}
    (s-P)^{-1}=\frac{s-1+P}{s(s-1)}\bigl(\Id+\epsilon E'_s\bigr),\quad s\in\Gamma.
  \end{equation*}
  (For compact fibres the existence statement is
  \cref{semiclassical:thm:projection-smoothing}, proved there by a Newton
  iteration, and the uniqueness up to homotopy is
  \cref{semiclassical:prop:projection-homotopy}; we give a direct argument.)
  Let $\mathfrak E=\frac1{2\pi i}\oint_\Gamma(s-P)^{-1}ds$ be the Riesz idempotent
  of $P$ for the part of the spectrum inside $\Gamma$; that it is an idempotent is
  the standard consequence of the resolvent identity. The residue of
  $(s-1+P)/(s(s-1))$ at $s=1$, the only pole inside $\Gamma$, is $P$, so
  \begin{equation*}
    \mathfrak E=P+\frac{\epsilon}{2\pi i}\oint_\Gamma\frac{s-1+P}{s(s-1)}
    E'_s\,ds\in p_0+Q_0+\epsilon\Psisl{-\infty},
  \end{equation*}
  which has the required properties (the kernels of $E'_s$ are supported in the
  same compact set as those of $E$).

  If $\mathfrak E'$ is a second quantization, then $P_t=(1-t)\mathfrak E+t
  \mathfrak E'$ is of the same form as $P$ above, uniformly in $t\in[0,1]$, so
  for small $\epsilon$ its Riesz idempotents $\mathfrak E_t$ are defined and
  depend smoothly on $t$; $\mathfrak E_0=\mathfrak E$ and $\mathfrak E_1=
  \mathfrak E'$, since the Riesz idempotent of an idempotent is itself. Parameters
  are handled in the same way.
\end{proof}

\subsection{Definition of the semiclassical index}
Let $x\in\Kc(T^*(M/B))$ be represented by a triple $(E_+,E_-,a)$, where $E_\pm$ are
bundles over $M$ and $a$ is a smooth section of $\Hom(E_+,E_-)$ over $T^*(M/B)$,
invertible outside a compact set $K_a$. Choose a smooth section $b$ of
$\Hom(E_-,E_+)$ such that
\begin{equation}\label{scindex:eq:b}
  r_+=\Id-ba\quad\text{and}\quad r_-=\Id-ab\quad\text{have compact support};
\end{equation}
for example $b=(1-\chi)a^{-1}$ with $\chi\in\CIc(T^*(M/B))$ equal to $1$ near
$K_a$. On $E=E_-\oplus E_+$ let $p_0$ be the projection onto $E_-$ and
$e_{a,b}$ the idempotent \eqref{scindex:eq:milnor}, formed pointwise. Then
$e_{a,b}-p_0$ is compactly supported.

\begin{definition}\label{scindex:def:Indsl}
  The \emph{semiclassical index} of $x$ is
  \begin{equation*}
    \Indsl{}(x)=d\bigl(\mathfrak E(\epsilon),p_0\bigr)\in K^0(B),
  \end{equation*}
  where $\mathfrak E$ is a semiclassical quantization of $e_{a,b}$ and
  $\epsilon>0$ is sufficiently small.
\end{definition}

\begin{theorem}\label{scindex:thm:Indsl}
  $\Indsl{}$ is a well-defined homomorphism $\Kc(T^*(M/B))\to K^0(B)$, and it has
  the following properties.
  \begin{enumerate}[label=(\roman*)]
  \item (Excision) If $U\subset M$ is open (a fibration over $B$) and
    $j:T^*(U/B)\hookrightarrow T^*(M/B)$ is the inclusion, then
    $\Indsl{M/B}\circ j_*=\Indsl{U/B}$.
  \item (Naturality) If $F:M\to M'$ is a diffeomorphism of fibrations over $B$,
    then $\Indsl{M'/B}\circ F_*=\Indsl{M/B}$. If $f:B'\to B$ is smooth,
    then $\Indsl{}(f^*x)=f^*\Indsl{}(x)$ for the pulled back fibration.
  \item (Degenerate fibration) For $\phi=\Id:B\to B$, the map
    $\Indsl{}$ on $\Kc(T^*(B/B))=K^0(B)$ is the identity.
  \item (Module property) $\Indsl{}(x\cdot\phi^*y)=\Indsl{}(x)\,y$ for $y\in K^0(B)$,
    where $\phi^*y$ is pulled back to $T^*(M/B)$.
  \end{enumerate}
\end{theorem}
\begin{proof}
  \emph{Independence of the choices.} By \cref{scindex:prop:quantize} and
  \ref{scindex:d-homotopy}, $d(\mathfrak E(\epsilon),p_0)$ does not depend on the
  quantization or on $\epsilon\le\epsilon_0$. If $b_0$, $b_1$ both satisfy
  \eqref{scindex:eq:b}, so does $b_t=(1-t)b_0+tb_1$, with supports in a fixed
  compact set; the quantizations of $e_{a,b_t}$ (\cref{scindex:prop:quantize}
  with parameter $t$) give a homotopy, so $\Indsl{}$ does not depend on $b$.

  \emph{Independence of the representative.} The equivalence relation on triples
  is generated by the following moves, which we check in turn. (b) Changing $a$
  on a compact set: if $a'=a$ outside a compact set, then $a_t=(1-t)a+ta'$ is
  invertible outside a fixed compact set, and the same $b$ satisfies
  \eqref{scindex:eq:b} for all $t$. (c) A homotopy $a_t$, invertible outside a
  fixed compact set: take $b_t=(1-\chi)a_t^{-1}$ with $\chi$ fixed. (a) Bundle
  isomorphisms $\phi_\pm$ of the lifted bundles: the dilation $\phi_\pm(t\zeta)$,
  $t\in[0,1]$, joins them to the lifts of their restrictions $\phi^0_\pm$ to the
  zero section, and along it the triple changes by a homotopy (c). For bundle
  isomorphisms $\phi_\pm^0$ over $M$, put $\Phi=\diag(\phi^0_-,\phi^0_+)$; then
  $e_{\phi^0_-a(\phi^0_+)^{-1},\phi^0_+b(\phi^0_-)^{-1}}=\Phi e_{a,b}\Phi^{-1}$,
  $\Phi\mathfrak E\Phi^{-1}$ is a quantization of it, $\Phi$ commutes with $p_0$,
  and \ref{scindex:d-conj} applies. (The same argument shows that the choice of
  the bundles $E_\pm$ over $M$ in a triple on $T^*(M/B)$ is immaterial.)
  (d) Stabilization by $(H,H,\Id)$: with $b=\Id$ on $H$, the extra block of
  $e_{a,b}$ is $p_0$, so by \ref{scindex:d-sum} nothing changes. Direct sums of
  triples give direct sums of idempotents, so $\Indsl{}$ is additive.

  (i) As in Step 1 of the proof of \cref{families:thm:excision},
  $x\in\Kc(T^*(U/B))$ has a representative $(E_+,E_-,a)$ with $E_\pm$ bundles over
  $U$, an isomorphism $\tau:E_+\to E_-$ over $U\setminus C$, $C\subset U$ compact,
  with both bundles trivialized there and $\tau=\Id$ in these trivializations, and
  $a=\tau$ over $U\setminus C$; $j_*x$ is represented by the extended bundles
  $\tilde E_\pm$ over $M$ and the extension $\tilde a$ of $a$ by $\Id$. The set
  where $a$ is not invertible is compact in $T^*(U/B)$; choose $\chi\in
  \CIc(T^*(U/B))$ equal to $1$ near it and $b=(1-\chi)\tilde a^{-1}$. Then
  $e_{\tilde a,b}=p_0$ outside $\supp\chi$. By \ref{scindex:S-symbol} a
  quantization on $U/B$ can be chosen with kernels supported in $K\times_BK$,
  $K\subset U$ compact; extended by $p_0$ it is a quantization on $M/B$. Now
  \ref{scindex:d-excision} gives the claim.

  (ii) The first statement follows from \ref{scindex:S-invariance} and
  \cref{scindex:lem:class}(iv): the push-forward of a quantization is a
  quantization of the push-forward symbol, and the loops are conjugated by a
  unitary. For the second, pull back the quantization and use the naturality of
  $\indT$.

  (iii) The fibres are points, $T^*(B/B)=B$, $\Psisl{-\infty}(B/B;E)$ consists
  of families of bundle maps, and the quantization of a symbol is the symbol
  itself. Since $B$ is compact, any $a$ and $b$ satisfy \eqref{scindex:eq:b};
  take $a=0$, $b=0$, so $r_\pm=\Id$ and $e_{0,0}=\diag(0,\Id_{E_+})$. By
  \ref{scindex:d-sum} and \ref{scindex:d-finite} (with $\Pi=\Id_{E_-}$ in the
  first block and $\Pi=\Id_{E_+}$ in the second),
  \begin{equation*}
    \Indsl{}([E_+]-[E_-])=d(0,\Id_{E_-})+d(\Id_{E_+},0)=-[E_-]+[E_+].
  \end{equation*}

  (iv) By additivity it suffices to take $y=[F]$, $F$ a bundle over $B$. The
  product $x\cdot\phi^*[F]$ is represented by $(E_+\otimes F,E_-\otimes F,
  a\otimes\Id_F)$ (the Clifford product with the triple $(F,0,0)$), so we may take
  $b\otimes\Id_F$, $e_{a,b}\otimes\Id_F$ and the quantization $\mathfrak E\otimes
  \Id_F$, which is in the calculus since $F$ is constant on the fibres. Now
  apply \ref{scindex:d-module}.
\end{proof}

\begin{remark}[Comparison with the quantization maps of \cref{ch:bott}]\label{scindex:rem:compare-bott}
  For the trivial fibration $\R^n\times X\to X$, $X$ compact, $T^*((\R^n\times X)/X)=
  \R^{2n}\times X$ and $\Indsl{(\R^n\times X)/X}$ coincides with the even
  isotropic--semiclassical quantization map $\Indsl{even}:\Kc^0(\R^{2n}\times X)\to
  K^0(X)$ of \cref{bott:prop:even-sl}. Indeed, both quantize projection-valued (or
  idempotent-valued) symbols equal to a constant $p_\infty$ outside a compact set;
  the quantizations differ only in the use of the Weyl instead of the left
  quantization and of the isotropic rescaling instead of the semiclassical one,
  and \cref{bott:def:sl-quantization} and the remark following it, together with
  \ref{scindex:S-euclid}, show that the two families of idempotents have the same
  semiclassical symbol and are joined by a homotopy of idempotent perturbations of
  $p_\infty$ for small $\epsilon$. Finally the relative class $\delta(Q)$ of
  \cref{bott:lem:relative-class} equals $d(Q,p_\infty)$: both are invariant under
  homotopies, and on the finite rank model of \cref{bott:lem:relative-class}(i)
  both are given by the range of the finite rank part, by
  \ref{scindex:d-finite}. (Projection data $p$ give idempotents $e$ through
  \cref{bott:lem:proj-model}(ii), whose construction is the Clifford rotation
  used in \cref{scindex:lem:milnor}; the classes agree by \ref{scindex:d-homotopy}.)
  Thus $\Indsl{}$ extends $\Indsl{even}$ from trivial bundles of phase planes to
  arbitrary fibrations, and its odd counterpart $\Indsl{odd}$ of
  \cref{scindex:sec:odd} extends the map of \cref{bott:prop:odd-sl}.
\end{remark}

\begin{remark}\label{scindex:rem:excision-trivial}
  Excision, which required some care for the analytic index (\cref{families:thm:excision}),
  is immediate here because nothing outside the support of the symbol enters the
  construction: for $\epsilon$ small the operators are $p_0$ plus families with
  kernels in an arbitrarily small neighbourhood of the fibre diagonal over the
  support. Note also that the fibres need not be compact.
\end{remark}

\section{The odd semiclassical index}\label{scindex:sec:odd}

In Melrose's presentation the basic object is the odd index, defined by quantizing
compactly supported perturbations of the identity. Recall from \cref{bott:def:Kc}
that $\Kc^{-1}(X)$ is the group of homotopy classes of smooth maps
$a:X\to\GL(N,\C)$ (stably in $N$) with $a=\Id$ outside a compact set, through
such maps, and $\Kc^{-2}(X)$ the analogous group of pointed loops.

\begin{proposition}\label{scindex:prop:odd}
  Let $a\in\CI(T^*(M/B);\GL(N,\C))$ with $a-\Id$ compactly supported, and let
  $Q\in\Psisl{-\infty}(M/B;\C^N)$ have $\sigsl(Q)=a-\Id$. For $\epsilon$ small,
  $\Id+Q(\epsilon)\in\Ginf_{\mathrm c}(M/B;\C^N)$, and its class
  \begin{equation*}
    \Indsl{odd}[a]=[\Id+Q(\epsilon)]\in K^{-1}(B)
  \end{equation*}
  depends only on $[a]\in\Kc^{-1}(T^*(M/B))$. This defines a homomorphism
  \begin{equation*}
    \Indsl{odd}:\Kc^{-1}(T^*(M/B))\to K^{-1}(B)
  \end{equation*}
  with the excision and naturality
  properties of \cref{scindex:thm:Indsl}. Applied to pointed loops (the fibration
  $M\times\Sphere\to B\times\Sphere$) it gives
  $\Indsl{odd,\Sphere}:\Kc^{-2}(T^*(M/B))\to K^{-2}(B)$.
\end{proposition}
\begin{proof}
  Invertibility for small $\epsilon$ is \cref{scindex:lem:invertible}. If $Q'$ is
  another choice, then $Q'-Q=\epsilon R$ with $R\in\Psisl{-\infty}$, and
  $\Id+Q+t\epsilon R=(\Id+Q)(\Id+t\epsilon(\Id+Q)^{-1}R)$ is invertible for
  $t\in[0,1]$ and $\epsilon$ small, by \cref{scindex:lem:uniform}; by
  \cref{scindex:lem:class} the class is unchanged. Homotopies of $a$ are
  quantized with a parameter, stabilization $a\oplus\Id$ adds an identity block,
  and $\Indsl{odd}[ab]=\Indsl{odd}[a]+\Indsl{odd}[b]$ because a quantization of $ab$ is
  $(\Id+Q_a)(\Id+Q_b)-\Id$. Excision and naturality are proved as for the even
  index.
\end{proof}

The even index is the odd one applied to loops. For $x=[(E_+,E_-,a)]$ as above,
the symbol-level loop
\begin{equation}\label{scindex:eq:symbol-loop}
  \ell_x(z)=\bigl(e_{a,b}z^{-1}+\Id-e_{a,b}\bigr)\bigl(p_0z+\Id-p_0\bigr)
\end{equation}
is a pointed loop of invertible maps which is the identity outside a compact set,
by the symbol-level version of \eqref{scindex:eq:loop-expansion}. (After
embedding $E$ in a trivial bundle and extending by $\Id$ it defines a class in
$\Kc^{-2}(T^*(M/B))$; this is the inverse of the Toeplitz index for compactly
supported classes.)

\begin{proposition}\label{scindex:prop:even-odd}
  $\Indsl{}(x)=\indT\bigl(\Indsl{odd,\Sphere}[\ell_x]\bigr)$.
\end{proposition}
\begin{proof}
  If $\mathfrak E=p_0+Q$ is a quantization of $e_{a,b}$, then by
  \eqref{scindex:eq:loop-expansion} $g_{\mathfrak E,p_0}-\Id\in\Psisl{-\infty}$
  with symbol $\ell_x-\Id$, so it is a valid quantization in
  \cref{scindex:prop:odd}. Hence $[g_{\mathfrak E,p_0}]=\Indsl{odd,\Sphere}[\ell_x]$, and
  $\Indsl{}(x)=\indT[g_{\mathfrak E,p_0}]$ by definition.
\end{proof}

\begin{remark}\label{scindex:rem:suspension}
  Melrose proves the compatibility of the odd and even semiclassical index maps
  with Bott periodicity (his Proposition~12.7, the commutativity of a square with
  vertical maps given by isotropic quantization) by introducing a joint
  isotropic--semiclassical calculus. With the definitions adopted here the even
  index is \emph{defined} through loops, and \cref{scindex:prop:even-odd} replaces
  that square; the isotropic quantization enters only in
  \cref{scindex:sec:thom}.
\end{remark}

\section{The analytic and semiclassical indices agree}\label{scindex:sec:analytic}

In this section $\phi:M\to B$ is a fibration of compact manifolds. We first
express the analytic index bundle as a relative index of idempotents, using
\cref{scindex:lem:milnor} at the operator level.

\begin{proposition}\label{scindex:prop:index-relative}
  Let $A\in\Psido0(M/B;E_+,E_-)$ be elliptic and let $B\in\Psido0(M/B;E_-,E_+)$
  satisfy $R_+=\Id-BA\in\Psido{-\infty}(M/B;E_+)$ and $R_-=\Id-AB\in
  \Psido{-\infty}(M/B;E_-)$. Then $e_{A,B}$ of \eqref{scindex:eq:milnor} is a
  compactly supported idempotent perturbation of $p_0$ on $E_-\oplus E_+$, and
  \begin{equation*}
    d(e_{A,B},p_0)=\inda\bigl([\sigma_0(A)]\bigr).
  \end{equation*}
  The same holds for the operators of \cref{scindex:rem:adjoin}, i.e.\ with
  $\C^L$ adjoined to $E_+$ or $E_-$.
\end{proposition}
\begin{proof}
  The entries of $e_{A,B}-p_0$ are products of elements of $\Psido0(M/B)$ with at
  least one factor $R_\pm$, so they are smoothing families; $e_{A,B}$ is an
  idempotent by \cref{scindex:lem:milnor}. Write $I(A,B)=d(e_{A,B},p_0)$.

  \emph{Step 1: $I$ depends only on $A$, and is homotopy invariant.} If $B'$ is a
  second such parametrix, then $B_t=(1-t)B+tB'$ has $\Id-B_tA=(1-t)R_++tR'_+$
  and similarly on the other side, so $e_{A,B_t}$ is a smooth family of compactly
  supported idempotent perturbations and \ref{scindex:d-homotopy} applies. If
  $(A_t,B_t)$ is a smooth family of such pairs, the same argument applies. We
  write $I(A)$.

  \emph{Step 2: stabilization.} Let $S:\C^L\to\CI(M;E_-)$ be given by $L$ smooth
  sections, such that $A\oplus S:L^2(Z_b;E_+)\oplus\C^L\to L^2(Z_b;E_-)$ is
  surjective for every $b$. Consider the operators $A_t=(A,tS)$ from $E_+\oplus\C^L$
  to $E_-$, in the sense of \cref{scindex:rem:adjoin}, with the fixed
  $B'=\bigl(\begin{smallmatrix}B\\0\end{smallmatrix}\bigr)$. Then
  \begin{equation*}
    \Id-B'A_t=\begin{pmatrix}R_+&-tBS\\0&\Id_{\C^L}\end{pmatrix},\qquad
    \Id-A_tB'=R_-,
  \end{equation*}
  both smoothing in the extended sense, so $I(A_0)=I(A_1)$. For $t=0$, $A_0$ is
  the direct sum of $A$ and the zero map $\C^L\to\{0\}$, and $B'$ is the direct
  sum of $B$ and $0$; for the second summand $R_+=\Id_{\C^L}$, $R_-=0$ and the
  idempotent is $\Id_{\C^L}$, with $p_0=0$ there. By \ref{scindex:d-sum},
  \ref{scindex:d-conj} (to reorder the summands) and \ref{scindex:d-finite},
  \begin{equation*}
    I(A,S)=I(A_1)=I(A_0)=I(A)+[\C^L].
  \end{equation*}

  \emph{Step 3: the surjective case.} Write $A'=(A,S)$ and let
  $\mathcal N=\mathcal N(A,S)$ be the null bundle, a smooth bundle of rank $r$
  whose fibres consist of smooth sections (\ref{scindex:K-inda}); the orthogonal
  projections $\Pi_{\mathcal N}$ onto its fibres form a smooth family of finite
  rank smoothing operators, being given locally by orthonormal frames. The
  family $C=A'^*A'+\Pi_{\mathcal N}$ is elliptic in the extended sense, and
  injective, since $\inner{Cu}{u}=\norm{A'u}^2+\norm{\Pi_{\mathcal N}u}^2$; being
  self-adjoint and Fredholm of index $0$, it is invertible. Its inverse is a smooth
  family in the extended calculus: in a local trivialization $C$ is norm-smooth
  (\cref{product:lem:norm-continuity}), hence so is $V=C^{-1}$, and if $Q$ is a parametrix with
  $QC=\Id-R$, $CQ=\Id-R'$, then $V=Q+RQ+RVR'$, where $RVR'$ has a smooth kernel
  depending smoothly on $b$ (as in the proof of \cref{bott:lem:inverse}). Put $G=C^{-1}A'^*$. Since
  $C\Pi_{\mathcal N}=\Pi_{\mathcal N}$ we have $C^{-1}\Pi_{\mathcal N}=
  \Pi_{\mathcal N}=\Pi_{\mathcal N}C^{-1}$, hence
  \begin{equation*}
    GA'=C^{-1}(C-\Pi_{\mathcal N})=\Id-\Pi_{\mathcal N},\qquad
    A'GA'=A'(\Id-\Pi_{\mathcal N})=A',
  \end{equation*}
  and, $A'$ being surjective, $A'G=\Id$. So $G$ is a parametrix of $A'$ with
  $R_+=\Pi_{\mathcal N}$, $R_-=0$. Since $A'\Pi_{\mathcal N}=0$ and
  $\Pi_{\mathcal N}G=\Pi_{\mathcal N}C^{-1}A'^*=(A'\Pi_{\mathcal N})^*=0$,
  \eqref{scindex:eq:milnor} gives
  \begin{equation*}
    e_{A',G}=\begin{pmatrix}\Id&0\\0&\Pi_{\mathcal N}\end{pmatrix}=p_0+
    \Pi,\qquad\Pi=\diag(0,\Pi_{\mathcal N}),
  \end{equation*}
  with $p_0\Pi=\Pi p_0=0$. By \ref{scindex:d-finite}, $I(A,S)=[\mathcal N]$.

  Combining the steps, $I(A)=[\mathcal N(A,S)]-[\C^L]=\inda[\sigma_0(A)]$.
\end{proof}

\begin{remark}\label{scindex:rem:index-trace}
  For $B$ a point, the rank of $d(e_{A,B},p_0)$ is $\Tr(e_{A,B}-p_0)=
  \Tr(R_+^2)-\Tr(R_-^2)$, which is the formula $\ind A=\Tr(R_+^k)-\Tr(R_-^k)$ of
  \cref{index:thm:trace-formula} for $k=2$. Indeed, $\Tr$ of the relative index of two
  idempotents differing by a trace class operator is the difference of their
  ranks.
\end{remark}

Next we need parametrices in the semiclassical calculus.

\begin{lemma}\label{scindex:lem:sl-parametrix}
  Let $a\in\Symph0(T^*(M/B);\Hom(E_+,E_-))$ be invertible outside a compact set,
  and $b\in\Symph0(T^*(M/B);\Hom(E_-,E_+))$ satisfy \eqref{scindex:eq:b}. Let
  $A\in\Psisl0(M/B;E_+,E_-)$ be a quantization of $a$ with symbols independent of
  $\epsilon$, as in \ref{scindex:S-order0}. Then there is
  $B\in\Psisl0(M/B;E_-,E_+)$ with $\sigsl(B)=b$ and
  \begin{equation*}
    \Id-BA\in\Psisl{-\infty}(M/B;E_+),\qquad\Id-AB\in\Psisl{-\infty}(M/B;E_-).
  \end{equation*}
\end{lemma}
\begin{proof}
  Since $ba=\Id$ near infinity, the boundary values of $a$ and $b$ are inverse
  to each other. Let $B_0$ be a quantization of $b$ with symbols independent of
  $\epsilon$; by \ref{scindex:S-order0}, $\sigma_0(B_0(\epsilon))=
  \sigma_0(A(\epsilon))^{-1}$ for all $\epsilon\in[0,1]$. Then
  $R=\Id-B_0A\in\Psisl0$ has vanishing principal symbols, so $R\in\Psisl{-1}$,
  and $\sigsl(R)=r_+$ is compactly supported. Choose $R_1\in\Psisl{-\infty}$ with
  $\sigsl(R_1)=r_+$. Then $R-R_1\in\Psisl{-1}$ has vanishing semiclassical symbol,
  so $R-R_1=\epsilon T$ with $T\in\Psisl{-1}$. Let $F\in\Psisl0$ be an asymptotic
  sum of $\sum_{k\ge0}\epsilon^kT^k$. Then $F(\Id-\epsilon T)-\Id\in
  \bigcap_K\epsilon^K\Psisl{-K}=\Psiinf$, $\sigsl(F)=\Id$, and with $B_1=FB_0$,
  \begin{equation*}
    \Id-B_1A=\Id-F(\Id-R_1-\epsilon T)=\bigl(\Id-F(\Id-\epsilon T)\bigr)+FR_1
    \in\Psisl{-\infty}.
  \end{equation*}
  In the same way there is $B_2$ with $\sigsl(B_2)=b$ and $\Id-AB_2\in
  \Psisl{-\infty}$. Then $B_1-B_2=B_1(\Id-AB_2)-(\Id-B_1A)B_2\in\Psisl{-\infty}$,
  so $B=B_1$ satisfies $\Id-AB=(\Id-AB_2)-A(B_1-B_2)\in\Psisl{-\infty}$.
\end{proof}

\begin{theorem}\label{scindex:thm:analytic=sl}
  For every fibration of compact manifolds $\phi:M\to B$,
  \begin{equation*}
    \Indsl{}=\inda:\Kc(T^*(M/B))\longrightarrow K^0(B).
  \end{equation*}
\end{theorem}
\begin{proof}
  Let $x\in\Kc(T^*(M/B))$. As in \cref{ch:bott,ch:families} we may represent $x$
  by $(E_+,E_-,a)$ with $E_\pm$ bundles over $M$ and $a$ homogeneous of degree $0$
  and invertible for $\abs\zeta\ge1$ (for some metric), extended smoothly to
  $T^*(M/B)$; then $a\in\Symph0$. Let $b=(1-\chi(\abs\zeta))a^{-1}$ with
  $\chi\in\CIc([0,\infty))$ equal to $1$ on $[0,1]$; it is a classical symbol of
  order $0$ satisfying \eqref{scindex:eq:b}. Choose $A\in\Psisl0(M/B;E_+,E_-)$
  quantizing $a$ with symbols independent of $\epsilon$, and $B$ as in \cref{scindex:lem:sl-parametrix}.

  For each fixed $\epsilon\in(0,1]$, $A(\epsilon)$ is an elliptic family in
  $\Psido0(M/B;E_+,E_-)$ with $\sigma_0(A(\epsilon))=a|_{\abs\zeta\ge1}$, so
  $[\sigma_0(A(\epsilon))]=x$, and $B(\epsilon)$ is a parametrix modulo smoothing
  families. By \cref{scindex:prop:index-relative},
  \begin{equation*}
    d\bigl(e_{A(\epsilon),B(\epsilon)},p_0\bigr)=\inda(x)\quad\text{for every
    }\epsilon\in(0,1].
  \end{equation*}
  On the other hand, $e_{A,B}-p_0$ is a matrix of products of elements of
  $\Psisl0$, each containing a factor $R_\pm\in\Psisl{-\infty}$, so
  $e_{A,B}-p_0\in\Psisl{-\infty}(M/B;E_-\oplus E_+)$; by multiplicativity of
  $\sigsl$ its symbol is $e_{a,b}-p_0$, where $e_{a,b}$ is formed from
  $a,b,r_\pm=\sigsl(R_\pm)$; and $e_{A,B}$ is an idempotent for every $\epsilon$.
  So $e_{A,B}$ is a semiclassical quantization of $e_{a,b}$ in the sense of
  \cref{scindex:prop:quantize}, and by definition $\Indsl{}(x)=
  d(e_{A(\epsilon),B(\epsilon)},p_0)$ for small $\epsilon$.
\end{proof}

\begin{corollary}[Excision for the analytic index]\label{scindex:cor:excision}
  In the situation of \cref{families:thm:excision} (two fibrations
  of compact manifolds $M_1$, $M_2$ over $B$ containing a common open subset $U$),
  $\inda\circ(i_1)_*=\inda\circ(i_2)_*$ on $\Kc(T^*(U/B))$.
\end{corollary}
\begin{proof}
  By \cref{scindex:thm:analytic=sl} and \cref{scindex:thm:Indsl}(i),(ii),
  $\inda\circ(i_j)_*=\Indsl{M_j/B}\circ(i_j)_*=\Indsl{U/B}$ for $j=1,2$.
\end{proof}

\begin{remark}\label{scindex:rem:compare-melrose-12.1}
  Melrose's proof of \cref{scindex:thm:analytic=sl} (his Theorem~12.1) is
  organized around a common generalization of the two index maps: triples
  $(\pi^+,\pi^-,a)$ of families of projections on the radially compactified
  cotangent bundle with an isomorphism $a$ of their ranges at infinity, quantized
  as Toeplitz-type operators $\Pi^-A'\Pi^+$ with $\Pi^-$ a semiclassical
  quantization of $\pi^-$. The analytic index corresponds to constant $\pi^\pm$,
  the semiclassical index to $a=\Id$, and the equality follows from homotopy
  invariance of the general index and the fact that both kinds of data are
  retracts of the general one. The idempotent $e_{A,B}$ plays the role of these
  triples here, and reduces the argument to the multiplicativity of the
  semiclassical symbol.
\end{remark}

\section{The Thom class and the harmonic oscillator}\label{scindex:sec:thom}

Let $\varpi:\mathcal V\to M$ be a real vector bundle of rank $k$ with a metric,
over a fibration of compact manifolds $M\to B$ (or over $M$ itself, which is the
case $B=M$). The fibration $\mathcal V\to M$ has non-compact fibres; it is an
open subset of the sphere bundle $\tilde M=S(\mathcal V\oplus\R)$ via inverse
stereographic projection $k_{\mathcal V}$, so its semiclassical index is defined.

\begin{theorem}\label{scindex:thm:thom}
  The Thom class has semiclassical index one:
  \begin{equation*}
    \Indsl{\mathcal V/M}(\lambda_{\mathcal V})=1=[\C]\in K^0(M).
  \end{equation*}
\end{theorem}

The proof occupies the rest of this subsection. We first describe the model on a
single fibre, then compare the semiclassical quantization with the harmonic
oscillator.

\subsection{The model}
On $\R^k$, with variables $X=(w,\xi)\in\R^{2k}$, let $\Lambda=\Lambda\C^k=
\Lambda^{\mathrm{ev}}\oplus\Lambda^{\mathrm{odd}}$ with its standard Hermitian
structure, $\{e_j\}$ the standard basis, $e_j\wedge$ exterior multiplication and
$\iota_j$ its adjoint. Let
\begin{equation*}
  c(X)=\sum_j(w_j+i\xi_j)\,e_j\wedge-\sum_j(w_j-i\xi_j)\,\iota_j
  :\Lambda^{\mathrm{ev}}\longrightarrow\Lambda^{\mathrm{odd}},
\end{equation*}
which is Clifford multiplication $c(u)=u\wedge\cdot-\iota_u$ by $u=w+i\xi$ as in
\cref{ch:families}; $c(X)^*=-c(X)$ and $c(X)^*c(X)=\abs X^2$. The group $O(k)$
acts on $\R^{2k}$ by $g(w,\xi)=(gw,g\xi)$ (the cotangent lift of its action on
$\R^k$) and on $\Lambda$ by $\Lambda g$, and $c(gX)=\Lambda g\,c(X)\,
(\Lambda g)^{-1}$.

Let $C_j=D_{w_j}+iw_j$, $A_j=D_{w_j}-iw_j$ be the creation and annihilation
operators of \cref{isotropic:sec:harmonic}, and $D=\sum_jA_j\,e_j\wedge$ the Koszul
operator, $D^*=\sum_jC_j\iota_j$. The left symbol of $A_j$ is
$\xi_j-iw_j=-i(w_j+i\xi_j)$ and that of $C_j$ is $\xi_j+iw_j=i(w_j-i\xi_j)$; since
$c$ is a sum of terms each of which is linear in $w$ or in $\xi$, its quantization
involves no ordering ambiguity, and
\begin{equation}\label{scindex:eq:koszul}
  \operatorname{Op}_1(c)=i(D+D^*)\big|_{\Lambda^{\mathrm{ev}}}=:i\mathcal D.
\end{equation}
Using $[A_j,C_l]=2\delta_{jl}$, $H=\sum_jC_jA_j+k$ and $e_j\wedge\iota_l+\iota_l
e_j\wedge=\delta_{jl}$,
\begin{equation*}
  (D+D^*)^2=DD^*+D^*D=\sum_{j,l}C_lA_j(e_j\wedge\iota_l+\iota_le_j\wedge)
  +2\sum_je_j\wedge\iota_j=H-k+2\deg
\end{equation*}
on $\Lambda$, since $D^2=0$ (compare \cref{isotropic:ex:koszul}). By \cref{isotropic:prop:H-inverse,isotropic:prop:HO-spectrum}, $H\ge k$ with
$\Nul(H-k)=\C u_0$, $u_0=\pi^{-k/4}e^{-\abs w^2/2}$. Hence
$\Nul(\mathcal D)=\Nul(D+D^*)=\C u_0\subset\Lambda^0$, and on
$\Lambda^{\mathrm{odd}}$, $\mathcal D\mathcal D^*=H-k+2\deg\ge2$ is invertible,
with inverse in $\Psiiso{-2}(\R^k;\Lambda^{\mathrm{odd}})$
(\cref{isotropic:prop:gen-inverse}). Put $G_0=\mathcal D^*(\mathcal D\mathcal D^*)^{-1}\in
\Psiiso{-1}$. Then
\begin{equation}\label{scindex:eq:G0}
  \mathcal DG_0=\Id,\qquad G_0\mathcal D=\Id-\Pi_0,
\end{equation}
where $\Pi_0$ is the orthogonal projection onto $\C u_0$ ($G_0\mathcal D$ is the
self-adjoint idempotent with range $\Ran\mathcal D^*=(\Nul\mathcal D)^\perp$).
All these operators commute with the action $(g\cdot u)(w)=\Lambda g\,u(g^{-1}w)$
of $O(k)$, and $u_0$ is invariant.

Now let
\begin{equation*}
  a=c,\qquad b=-(1-\chi(\abs X))\frac{c(X)}{\abs X^2},\qquad
  r_\pm=\chi(\abs X)\Id,
\end{equation*}
with $\chi\in\CIc([0,\infty))$ equal to $1$ near $[0,1]$. Since $c^{-1}=c^*/\abs X^2
=-c/\abs X^2$, $b$ satisfies \eqref{scindex:eq:b} with these $r_\pm$, it is a
symbol of order $-1$, and $e_{a,b}$ is an $O(k)$-equivariant family of
idempotents on $\Lambda^{\mathrm{odd}}\oplus\Lambda^{\mathrm{ev}}$ equal to $p_0$
(the projection onto $\Lambda^{\mathrm{odd}}$) for $\abs X$ large.

\subsection{Transport to the bundle}
Let $P\to M$ be the orthonormal frame bundle of $\mathcal V$, so that a point
$f\in P_m$ is an isometry $f:\R^k\to\mathcal V_m$, inducing $\Lambda\C^k\cong
\Lambda W_m$ and $L^2(\R^k;\Lambda)\cong L^2(\mathcal V_m;\Lambda W_m)$. An
$O(k)$-equivariant operator $T$ on $L^2(\R^k;\Lambda)$ defines operators
$T_m=f_*Tf_*^{-1}$ on $L^2(\mathcal V_m;\Lambda W_m)$, independent of $f\in P_m$, and
using smooth local sections of $P$ one sees that $(T_m)_{m\in M}$ is a smooth
family whenever the kernel of $T$ is smooth (or Schwartz). The cotangent lift of $f$
identifies $\R^{2k}$ with $T^*_{(m,\cdot)}(\mathcal V_m)=\mathcal V_m\oplus
\mathcal V_m^*$, and under this identification $e_{a,b}$ becomes the idempotent
formed, as in \cref{scindex:def:Indsl}, from the triple $(\Lambda^{\mathrm{ev}}W,
\Lambda^{\mathrm{odd}}W,c)$ representing $\lambda_{\mathcal V}$.

Let $Q_0$ be the operator with kernel $\epsilon^{-k}\check q(w,(w-w')/\epsilon)
\psi(\abs{w'})$, $q=e_{a,b}-p_0$, where $\psi\in\CIc([0,\infty))$ is $1$ on a
large interval; it is $O(k)$-equivariant, and $\Ope(q)-Q_0$ has a kernel which is
$O(\epsilon^\infty)$ in $\cS$. Running the construction of
\cref{scindex:prop:quantize} with this $Q_0$ gives an $O(k)$-equivariant
semiclassical quantization $\mathfrak F=p_0+Q$ of $e_{a,b}$ on $\R^k$ (the Riesz
integral preserves equivariance). Transported, $\mathfrak F_m$ is a semiclassical
quantization of the idempotent of $\lambda_{\mathcal V}$, by
\ref{scindex:S-invariance}. Hence
\begin{equation}\label{scindex:eq:thom-F}
  \Indsl{\mathcal V/M}(\lambda_{\mathcal V})=d(\mathfrak F_{\mathcal V},p_0),
\end{equation}
where the subscript denotes the transported family.

\subsection{Comparison with the harmonic oscillator}
Let $U_\epsilon u(w)=\epsilon^{k/4}u(\epsilon^{1/2}w)$, a unitary operator on
$L^2(\R^k;\Lambda)$ commuting with $O(k)$. Changing variables in
\eqref{scindex:eq:Ope},
\begin{equation}\label{scindex:eq:dilation}
  U_\epsilon\Ope(q)U_\epsilon^{-1}=\operatorname{Op}_1(q_\epsilon),\qquad
  q_\epsilon(w,\xi)=q(\epsilon^{1/2}w,\epsilon^{1/2}\xi).
\end{equation}
In particular $U_\epsilon\Ope(c)U_\epsilon^{-1}=\epsilon^{1/2}\operatorname{Op}_1(c)
=i\epsilon^{1/2}\mathcal D$ by \eqref{scindex:eq:koszul}: the semiclassical
quantization of the Thom symbol is, up to a unitary rescaling and a factor, the
Koszul complex of the harmonic oscillator. Conjugation by $U_\epsilon$ maps
$\Psiiso{m}(\R^k)$ to itself, and an operator has a Schwartz kernel if and only if
its conjugate does. We write $\Psiiso{-\infty}$ for the operators with Schwartz
kernels, an ideal in $\bigcup_m\Psiiso m$ (\cref{isotropic:lem:residual-kernels,isotropic:thm:algebra}).

\begin{lemma}\label{scindex:lem:thom-parametrix}
  Let $A=\Ope(c)$ and $B_0=\Ope(b)$. For $\epsilon$ small there is
  $B\in U_\epsilon^{-1}\Psiiso{-1}U_\epsilon$, $O(k)$-equivariant, with
  $R_+=\Id-BA$ and $R_-=\Id-AB$ in $\Psiiso{-\infty}$ and
  \begin{equation*}
    \norm{e_{A,B}-p_0-\Ope(e_{a,b}-p_0)}_{L^2}\le C\epsilon.
  \end{equation*}
\end{lemma}
\begin{proof}
  Since $c$ is linear in $\xi$ and in $w$, the expansions of
  \ref{scindex:S-euclid} are exact:
  \begin{equation*}
    B_0A=\Ope\Bigl(ba+\epsilon\sum_j\pa_{\xi_j}b\,D_{w_j}c\Bigr)
    =\Id-\Ope(r_+)+\epsilon\Ope(s),\qquad s=-i\sum_j\pa_{\xi_j}b\,c(e_j),
  \end{equation*}
  where $c(e_j)=e_j\wedge-\iota_j$ is constant, and similarly
  $AB_0=\Id-\Ope(r_-)+\epsilon\Ope(s')$ with $s'=\sum_j\pa_{\xi_j}c\,D_{w_j}b$.
  Both $s,s'$ are symbols of order $-2$ in $X$, so they lie in $\Sym0(\R^k;\R^k)$
  and $T_1=\epsilon\Ope(s)$ has norm $\le C\epsilon$ by \ref{scindex:S-euclid}.
  For small $\epsilon$, $\Id+T_1$ is invertible on $L^2$; after conjugation by
  $U_\epsilon$ it is an elliptic element of $\Psiiso0$ (its symbol is
  $1+\epsilon s_\epsilon$, with $s_\epsilon\in\Gs{-2}$), so its inverse is in
  $\Psiiso0$ (an elliptic isotropic operator invertible on $L^2$ has trivial null
  spaces on $\cSp$, since these lie in $\cS$, and its generalized inverse in
  \cref{isotropic:prop:gen-inverse} is then an inverse). Set $B=(\Id+T_1)^{-1}B_0$. Then
  \begin{equation*}
    R_+=\Id-BA=(\Id+T_1)^{-1}\Ope(r_+)\in\Psiiso{-\infty},\qquad
    \norm{R_+-\Ope(r_+)}+\norm{B-B_0}\le C\epsilon.
  \end{equation*}
  Next, $R_-A=AR_+\in\Psiiso{-\infty}$, and $A$ has the exact right inverse
  $G=U_\epsilon^{-1}(i\epsilon^{1/2})^{-1}G_0U_\epsilon$ by \eqref{scindex:eq:G0},
  so $R_-=R_-AG=AR_+G\in\Psiiso{-\infty}$. For the norm, write
  \begin{equation*}
    R_-=\Ope(r_-)-\epsilon\Ope(s')+AT_1(\Id+T_1)^{-1}B_0 .
  \end{equation*}
  Here $AT_1=\epsilon\Ope(c\,\#_\epsilon s)$ and $c\#_\epsilon s=cs+\epsilon
  \sum_j\pa_{\xi_j}c\,D_{w_j}s$ (exact, $c$ being linear in $\xi$) is a symbol of
  order $-1$, so $\norm{AT_1}\le C\epsilon$ and $\norm{R_--\Ope(r_-)}\le C\epsilon$.

  Finally, the entries of $e_{A,B}-p_0$ are $-R_-^2$, $(\Id+R_-)AR_+$, $R_+B$ and
  $R_+^2$; all of them are in $\Psiiso{-\infty}$, since $R_\pm$ are. By the
  estimates just proved and \ref{scindex:S-euclid}, $R_-^2$, $R_+B$ and $R_+^2$
  are within $C\epsilon$ of $\Ope(r_-^2)$, $\Ope(r_+b)$, $\Ope(r_+^2)$. For the
  remaining entry, $AR_+=A\Ope(r_+)-AT_1(\Id+T_1)^{-1}\Ope(r_+)$ and
  $A\Ope(r_+)=\Ope(cr_++\epsilon\sum_j\pa_{\xi_j}c\,D_{w_j}r_+)$ exactly, with
  $r_+$ compactly supported; so $AR_+$ is within $C\epsilon$ of $\Ope(cr_+)$, and
  $(\Id+R_-)AR_+$ within $C\epsilon$ of $\Ope((1+r_-)cr_+)$. These symbols are the
  entries of $e_{a,b}-p_0$. Equivariance holds at each step.
\end{proof}

\begin{proof}[Proof of \cref{scindex:thm:thom}]
  By \cref{scindex:prop:quantize} and \cref{scindex:lem:uniform},
  $\norm{\mathfrak F-p_0-\Ope(e_{a,b}-p_0)}\le C\epsilon$, so with $e=e_{A,B}$ from
  \cref{scindex:lem:thom-parametrix}, $\norm{e-\mathfrak F}\le C\epsilon$. Both
  $e-p_0$ and $\mathfrak F-p_0$ have Schwartz kernels. For $t\in[0,1]$ put
  $e_t=(1-t)e+t\mathfrak F$; a short computation gives
  \begin{equation*}
    e_t^2-e_t=-t(1-t)(e-\mathfrak F)^2,
  \end{equation*}
  which has norm $\le C\epsilon^2$ and a Schwartz kernel. As in the proof of
  \cref{scindex:prop:quantize}, for small $\epsilon$ the Riesz idempotents
  $\mathfrak e_t=\frac1{2\pi i}\oint_{\abs{s-1}=1/2}(s-e_t)^{-1}ds$ are defined,
  $\mathfrak e_0=e$, $\mathfrak e_1=\mathfrak F$, and
  \begin{equation*}
    (s-e_t)^{-1}=\frac{s-1+e_t}{s(s-1)}\Bigl(\Id-\frac{e_t^2-e_t}{s(s-1)}
    \Bigr)^{-1}.
  \end{equation*}
  The inverse on the right is $\Id$ plus an operator with Schwartz kernel
  (\cref{isotropic:cor:invertible}: if $K$ has a Schwartz kernel and $\Id+K$ is invertible,
  $(\Id+K)^{-1}=\Id-K+K(\Id+K)^{-1}K$, and the last term again has a Schwartz kernel,
  depending smoothly on parameters). Hence $\mathfrak e_t-p_0$ has a Schwartz
  kernel depending smoothly on $t$, and $\mathfrak e_t$ is $O(k)$-equivariant.
  Transporting to $\mathcal V$ and using \cref{scindex:rem:schwartz} and
  \ref{scindex:d-homotopy},
  \begin{equation*}
    d(\mathfrak F_{\mathcal V},p_0)=d(e_{\mathcal V},p_0),\qquad e=e_{A,B}.
  \end{equation*}
  Now $G$ (the exact right inverse of $A$ above) satisfies $\Id-GA=\Pi$, the
  orthogonal projection onto $\C U_\epsilon^{-1}u_0$, and $\Id-AG=0$. The
  linear path $B_t=(1-t)B+tG$ has $\Id-B_tA=(1-t)R_++t\Pi$ and $\Id-AB_t=
  (1-t)R_-$, both with Schwartz kernels, so as in Step~1 of the proof of
  \cref{scindex:prop:index-relative}, $d(e_{\mathcal V},p_0)=
  d((e_{A,G})_{\mathcal V},p_0)$. Since $A\Pi=0$ and $\Pi G=0$ (as
  $\Pi\mathcal D^*=0$), \eqref{scindex:eq:milnor} gives $e_{A,G}=p_0+
  \diag(0,\Pi)$. The transported family $\Pi_{\mathcal V}$ is the projection onto
  the span of the section $m\mapsto(U_\epsilon^{-1}u_0)_m$, which is well defined
  because $u_0$ is $O(k)$-invariant and nowhere zero; so $\Ran\Pi_{\mathcal V}$ is
  the trivial line bundle and, by \ref{scindex:d-finite}, $d((e_{A,G})_{\mathcal V},
  p_0)=[\C]$. With \eqref{scindex:eq:thom-F} this proves the theorem.
\end{proof}

\begin{remark}\label{scindex:rem:why-trivial}
  The analogous statement for the analytic index, \ref{families:TH}, is
  traditionally proved by computing the index of the Bott operator on the sphere
  $\Sphere^k$ as a representation of $O(k)$ and then twisting by the frame bundle.
  Here the corresponding representation is visible directly: it is the span of the
  Gaussian ground state, on which $O(k)$ acts trivially. Since
  $\lambda_{\mathcal V}=\beta_{\mathcal V\otimes\C}$, \cref{scindex:thm:thom} is the
  statement that the fibrewise quantization map of
  \cref{bott:rem:thom-quantization} takes the Bott class of the complex bundle
  $\mathcal V\otimes\C=\mathcal V\oplus\mathcal V^*$ to $1$; for a complexified real
  bundle the structure group reduces to $O(k)\subset\mathrm U(k)$, which acts
  honestly (not projectively) on $L^2(\R^k)$, so the patching issue mentioned there
  does not arise. This is Melrose's
  observation that ``the bundle of ground states of the harmonic oscillators is
  canonically trivial''. The analysis needed is the comparison
  \cref{scindex:lem:thom-parametrix} between semiclassical quantization at scale
  $\epsilon$ and isotropic quantization after rescaling by $\epsilon^{1/2}$.
\end{remark}

\begin{corollary}[The index of the Thom class]\label{scindex:cor:TH}
  Let $\tilde M=S(\mathcal V\oplus\R)\to M$ and $k_{\mathcal V}:\mathcal V
  \hookrightarrow\tilde M$ be as in statement \ref{families:TH} of \cref{ch:families}. Then
  \begin{equation*}
    \inda^{\tilde M/M}\bigl((k_{\mathcal V})_*\lambda_{\mathcal V}\bigr)=1\in
    K^0(M).
  \end{equation*}
  Thus \ref{families:TH} holds, and the proof of \cref{families:thm:AS-K} given in
  \cref{ch:families} is complete.
\end{corollary}
\begin{proof}
  By \cref{scindex:thm:analytic=sl} for the fibration $\tilde M\to M$ of compact
  manifolds, excision (\cref{scindex:thm:Indsl}(i)) and \cref{scindex:thm:thom},
  $\inda((k_{\mathcal V})_*\lambda_{\mathcal V})=\Indsl{\tilde M/M}
  ((k_{\mathcal V})_*\lambda_{\mathcal V})=\Indsl{\mathcal V/M}
  (\lambda_{\mathcal V})=1$.
\end{proof}

\begin{corollary}[Semiclassical normalization]\label{scindex:cor:normalization}
  For every $N\ge0$ and every compact $B$,
  \begin{equation*}
    \Indsl{(I^N\times B)/B}\circ\beta_N=\Id_{K^0(B)},\qquad\text{hence}\qquad
    \inda\circ(k_N)_*\circ\beta_N=\Id_{K^0(B)},
  \end{equation*}
  where $k_N:T^*I^N\times B\hookrightarrow T^*(\Torus^N\times B/B)$.
\end{corollary}
\begin{proof}
  Apply \cref{scindex:thm:thom} to the trivial bundle $\mathcal V=\R^N\times B$
  over $M=B$ (the degenerate fibration $B\to B$): $\Indsl{(\R^N\times B)/B}
  (\lambda_{\R^N})=1$. By the module property \cref{scindex:thm:Indsl}(iv),
  $\Indsl{}(\lambda_{\R^N}\cdot y)=y$ for $y\in K^0(B)$, and $\lambda_{\R^N}\cdot y=
  \lambda_{\R^N}\boxtimes y$. Under a diffeomorphism $I\cong\R$,
  $\lambda_{\R^N}=b_1^{\boxtimes N}$ (\ref{scindex:K-families}), so by naturality
  $\Indsl{(I^N\times B)/B}(\beta_N(y))=y$. The second statement follows by
  excision ($I^N\times B\subset\Torus^N\times B$) and
  \cref{scindex:thm:analytic=sl}.
\end{proof}
For $N=1$ and with $b_1=b_{\mathrm{iso}}$ (\eqref{bott:eq:biso}; both are
represented by $x+i\xi$ up to homotopy), the first identity is also
\cref{bott:prop:sl-bott-one}, via \cref{scindex:rem:compare-bott}.
The second identity is \cref{families:prop:normalization}, which
was proved there by the multiplicativity of the analytic index on the torus; here
it is obtained from the harmonic oscillator alone.

\section{The families index theorem}\label{scindex:sec:AS}

We now give the second proof of the K-theoretic form of the families index
theorem of \cref{ch:families}. Recall the data defining $\indt$
(\ref{scindex:K-families}): an embedding $M\hookrightarrow\R^N\times B$ over
$B$ with normal bundle $\mathcal V\to M$, a tubular neighbourhood
$\Theta:\mathcal V\cong\Omega\subset I^N\times B$, the Thom homomorphism $\tau$,
the extension $j_*$ from $\Omega$ to $I^N\times B$, and the Bott isomorphism
$\beta_N$.

\begin{theorem}[Atiyah--Singer]\label{scindex:thm:AS}
  For every fibration of compact manifolds $\phi:M\to B$,
  \begin{equation*}
    \inda=\Indsl{}=\indt:\Kc(T^*(M/B))\longrightarrow K^0(B).
  \end{equation*}
  Moreover $\Indsl{(I^N\times B)/B}=\beta_N^{-1}$ and, for the fibration
  $\mathcal V\to B$,
  \begin{equation}\label{scindex:eq:thom-sl}
    \Indsl{\mathcal V/B}\circ\tau=\Indsl{M/B}.
  \end{equation}
\end{theorem}
\begin{proof}
  The first equality is \cref{scindex:thm:analytic=sl}.

  \emph{The trivial fibration.} By \cref{scindex:cor:normalization},
  $\Indsl{(I^N\times B)/B}\circ\beta_N=\Id$, and $\beta_N$ is an isomorphism by
  Bott periodicity (\cref{ch:bott,ch:families}); hence $\Indsl{(I^N\times B)/B}
  =\beta_N^{-1}$.

  \emph{The Thom step \eqref{scindex:eq:thom-sl}.} Let $\tilde M=S(\mathcal
  V\oplus\R)\xrightarrow{\tilde\phi}M$ and $\hat k:T^*(\mathcal V/B)
  \hookrightarrow T^*(\tilde M/B)$ the open embedding induced by $k_{\mathcal V}$.
  The class $b=(k_{\mathcal V})_*\lambda_{\mathcal V}\in\Kc(T^*(\tilde M/M))$ has
  $\inda(b)=1$ by \cref{scindex:cor:TH}. By \cref{families:cor:mult} and
  \cref{families:lem:natural} (with a splitting of $T^*(\tilde M/B)$ extending the one used to define
  $\tau$),
  \begin{equation*}
    \inda^{M/B}(x)=\inda^{\tilde M/B}\bigl(b\cdot\tilde\phi^*x\bigr)
    =\inda^{\tilde M/B}\bigl(\hat k_*\tau(x)\bigr).
  \end{equation*}
  By \cref{scindex:thm:analytic=sl} for $\tilde M\to B$ and $M\to B$, and excision
  for $\Indsl{}$, the outer terms are $\Indsl{M/B}(x)$ and
  $\Indsl{\mathcal V/B}(\tau(x))$.

  \emph{Conclusion.} By naturality and excision for $\Indsl{}$
  (\cref{scindex:thm:Indsl}),
  \begin{equation*}
    \Indsl{\mathcal V/B}(\tau x)=\Indsl{\Omega/B}(\Theta_*\tau x)=
    \Indsl{(I^N\times B)/B}(j_*\Theta_*\tau x)=\beta_N^{-1}(j_*\Theta_*\tau x)
    =\indt(x).
  \end{equation*}
  Together with \eqref{scindex:eq:thom-sl} this gives $\Indsl{}=\indt$.
\end{proof}

\begin{remark}[Comparison of the two proofs]\label{scindex:rem:compare}
  Both proofs pass through the same diagram, the embedding of $M$ into
  $I^N\times B$ followed by the Thom homomorphism. In \cref{families:thm:AS-K} the
  analytic index was carried along it using (1) multiplicativity, proved with the
  product-type calculus of \cref{ch:product}, (2) excision, proved by a direct
  comparison of stabilized null spaces, (3) the normalization on the torus,
  proved by multiplicativity, and (4) the index of the Thom class, (TH), which was
  imported. Here (2) and (3) are replaced by excision for $\Indsl{}$ (which is
  trivial) together with \cref{scindex:thm:analytic=sl}, and by
  \cref{scindex:cor:normalization}; and (4) is proved
  (\cref{scindex:cor:TH}). Only (1) is still taken from \cref{ch:families}. In
  Melrose's treatment (1) is replaced as well, by multiplicativity of the
  semiclassical index under iterated fibrations, proved by an adiabatic limit; we
  explain this in the next section. With it, \eqref{scindex:eq:thom-sl} is
  proved directly and the proof becomes entirely semiclassical:
  $\inda=\Indsl{M/B}=\Indsl{\mathcal V/B}\circ\tau=\beta_N^{-1}\circ j_*\circ
  \Theta_*\circ\tau$.
\end{remark}

\section{Multiplicativity and the adiabatic limit}\label{scindex:sec:adiabatic}

This section describes, without complete proofs, how the remaining ingredient,
\eqref{scindex:eq:thom-sl}, and more generally the multiplicativity of the
semiclassical index, can be obtained by semiclassical methods.

\subsection{Iterated fibrations}
Consider a tower $\tilde M\xrightarrow{\tilde\phi}M\xrightarrow{\phi}B$ of
fibrations, $\psi=\phi\circ\tilde\phi$. A class $\tilde x\in\Kc(T^*(\tilde M/B))$
can be quantized in two steps. First, in the fibres of $\tilde\phi$ only: this
requires the fibres of $\tilde\phi$ to carry, as parameters, the covectors in
$T^*(M/B)$, i.e.\ one considers the pulled back fibration
$\tilde\phi^*:{}^{\mathrm{sl}}T^*(M/B)\times_M\tilde M\to{}^{\mathrm{sl}}T^*(M/B)$,
and quantizes $\tilde x$ semiclassically along its fibres. This gives a class
$\Indsl{\tilde\phi^*}(\tilde x)\in\Kc(T^*(M/B))$ (compactly supported in the
parameters). Second, one applies $\Indsl{M/B}$. Melrose's multiplicativity
theorem (his Proposition~12.15) states that
\begin{equation}\label{scindex:eq:mult-sl}
  \Indsl{\psi}=\Indsl{\phi}\circ\Indsl{\tilde\phi^*}.
\end{equation}
For the tower $\mathcal V\to M\to B$ and $\tilde x=\lambda_{\mathcal V}\cdot
\varpi^*x$ the inner index is $\Indsl{\mathcal V/M}(\lambda_{\mathcal V})\cdot x=x$
by \cref{scindex:thm:thom} and the module property (applied with parameters in
$T^*(M/B)$), so \eqref{scindex:eq:mult-sl} contains \eqref{scindex:eq:thom-sl}.

\subsection{The adiabatic calculus}
To compare the two quantizations one uses two parameters: $\epsilon$, the
semiclassical parameter in the directions of $M/B$, and $\delta$, an adiabatic
parameter in the fibres of $\tilde\phi$. For a single fibration
$Y\to Z$ with fibre $\tilde Z$ (the case $B=\mathrm{pt}$), an \emph{adiabatic
family} is a family $A(\delta)$ of smoothing operators on $Y$ whose kernels vanish
to infinite order as $\delta\downarrow0$ away from the fibre diagonal
$Y\times_ZY$ and near it have the form
\begin{equation*}
  \delta^{-m}A'\Bigl(\delta,z,\frac{z-z'}\delta,\tilde z,\tilde z'\Bigr),\quad
  m=\dim Z,
\end{equation*}
in local product coordinates, with $A'$ smooth, compactly supported in $z$ and
Schwartz in $(z-z')/\delta$, and smooth in the fibre variables $\tilde z$,
$\tilde z'$. They form an algebra with a multiplicative symbol sequence
\begin{equation*}
  \delta\Psi^{-\infty}_{\mathrm{ad}}(Y;E)\hookrightarrow\Psi^{-\infty}_{\mathrm{ad}}
  (Y;E)\xrightarrow{\ \sigma_{\mathrm{ad}}\ }\cS\bigl({}^{\mathrm{sl}}T^*Z;
  \Psido{-\infty}(\pi^*Y/{}^{\mathrm{sl}}T^*Z;E)\bigr),
\end{equation*}
the adiabatic symbol being a Schwartz family, parametrized by the semiclassical
cotangent bundle of the base, of smoothing operators on the fibres $\tilde Z$.
If $a$ is an invertible-at-infinity adiabatic symbol, its adiabatic
quantizations are invertible for small $\delta$, and this defines an index
$\Kc^{-1}(T^*Z)\to K^{-1}(\mathrm{pt})$ with \emph{operator-valued} symbols; by
deforming $a$ to finite rank (\cref{scindex:ex:frame}) one sees that it agrees with
the semiclassical index (Melrose's Proposition~12.14). Combining the two
parameters gives the adiabatic--semiclassical algebra, whose kernels are of the
form $\delta^{-n-k}A(\delta,\epsilon,(\tilde z-\tilde z')/\epsilon,z,
(z-z')/(\epsilon\delta))$ near the fibre diagonal, with two symbol maps (at
$\epsilon=0$ and at $\delta=0$) and a joint symbol; the relevant cotangent
bundle ${}^{\mathrm{sl\,ad}}T^*\tilde M$ is spanned by $df/\epsilon$ and
$dg/(\epsilon\delta)$, $f\in\CI(\tilde M)$, $g\in\CI(Z)$, and restricts to
${}^{\mathrm{sl}}T^*\tilde M$, to the adiabatic cotangent bundle and to $T^*\tilde M$
on the three faces. The multiplicativity \eqref{scindex:eq:mult-sl} is then the
statement that a family in this algebra quantizing $\tilde x$ restricts, at the
two faces, to the two composite quantizations, and that the index is constant
along the family.

For the Thom tower the same idea can be made concrete in the language of this
chapter. Let $A$ be a semiclassical quantization of the symbol of $x$ on $M/B$,
lifted to $\mathcal V$ by a connection, and let $\mathcal D_\epsilon$ be the
transported, rescaled Koszul operator of \cref{scindex:sec:thom} on the fibres of
$\mathcal V$, whose null space is spanned by the transported Gaussian at scale
$\epsilon^{1/2}$ and which is surjective. This is the hypothesis (H) of the
\cref{families:sec:mult} (\eqref{families:eq:H}), and the Clifford tensor product
$\mathcal D_\epsilon\otimes A$ formed there has relative index equal to that of
$A$, by the same linear algebra. The remaining step is to show that the
idempotent of \cref{scindex:lem:milnor} built from $\mathcal D_\epsilon\otimes A$
is homotopic, through compactly supported (or Schwartz) idempotent perturbations
of $p_0$, to a semiclassical quantization on $\mathcal V/B$ of the idempotent of
$\lambda_{\mathcal V}\cdot\varpi^*x$. In the analytic proof this step required the
product-type calculus; here it requires the estimate that the commutators of the
lifted operators with $\mathcal D_\epsilon$ are $O(\epsilon^{1/2})$ relative to
$\mathcal D_\epsilon$, which is what the joint semiclassical--isotropic
calculus provides. We do not develop that calculus here.

\section{The Chern character of the index bundle}\label{scindex:sec:chern}

Since $\inda=\Indsl{}=\indt$, the cohomological form of the index theorem,
\begin{equation*}
  \Ch(\inda x)=(-1)^n\pi_*\bigl(\Ch(x)\wedge\pi_0^*\Td(T(M/B)\otimes\C)\bigr),
\end{equation*}
follows from the K-theoretic form exactly as in \cref{families:thm:AS-H}, with the same
inputs and the same conventions (multiplicativity of $\Ch$ for products,
compatibility with extension and pull-back, the integral of the Chern character of
$b_1$, and the Riemann--Roch identity for the Thom class). We add two remarks on
the semiclassical point of view.

\begin{remark}[The odd form]\label{scindex:rem:odd-chern}
  Melrose states the result for the odd index: for $a\in\Kc^{-1}(T^*(M/B))$,
  $\Ch_{\mathrm{odd}}(\Indsl{odd}(a))$ is the fibre integral of
  $\Ch_{\mathrm{odd}}(a)$ against the Todd class of the fibration. The two forms
  are equivalent. Indeed, by \cref{scindex:prop:even-odd}, $\Indsl{}(x)=
  \indT(\Indsl{odd,\Sphere}[\ell_x])$, and \cref{families:prop:toeplitz-chern} shows that the Chern character
  of a Toeplitz index is the integral over the circle of the odd Chern character
  of the loop; the same computation, now with compact supports, identifies the
  integral over the circle of the odd Chern character of $\ell_x$ with the Chern
  character of $x$ (we do not carry this out). The signs depend on the orientation
  conventions, which are fixed in \cref{ch:families}.
\end{remark}

\begin{remark}[The Todd class]\label{scindex:rem:todd}
  The derivation shows that $\Ch(\inda x)$ is the fibre integral of
  $\Ch(x)$ against a class $\mathcal T\in\Hev(M)$ obtained by integrating
  $\Ch(\lambda_{\mathcal V})$ over the fibres of $T^*(\mathcal V/M)$, where
  $\mathcal V$ is the normal bundle of an embedding. That $\mathcal T$ does not
  depend on the embedding follows from \cref{scindex:thm:AS} (for $B$ a point and
  suitable $x$) or directly: the class is stable under adding trivial bundles,
  since $\Ch(\lambda_{\R})$ integrates to a constant, and any two normal bundles
  are stably isomorphic, being stable inverses of $T(M/B)$. Its identification with
  the Todd class of $T(M/B)\otimes\C$ is the Riemann--Roch computation of Atiyah and
  Hirzebruch recalled in \cref{ch:families}.
\end{remark}

\begin{remark}[Dirac families and spectral sections]\label{scindex:rem:dirac}
  For families of Dirac operators the index formula can be refined to the level
  of differential forms, through the heat kernel and Bismut's superconnection;
  see \cite{Bismut-Cheeger} and \cref{ch:eta}, where eta forms appear in the
  index theorem for families with boundary. In that setting one also needs
  \emph{spectral sections} (smooth families of projections which differ from the
  positive spectral projections of a self-adjoint family by smoothing operators;
  \cite{Melrose-APS}): such a family exists if and only if the odd index of the
  boundary family vanishes in $K^{1}$ of the base, an obstruction which is the
  odd counterpart of the even index studied here. These topics are not pursued
  in this chapter.
\end{remark}

\section{Exercises}

\begin{exercise}[Strong contractibility of the unitary group]\label{scindex:ex:strong}
  Let $H=L^2([0,1])$ and $\mathrm U(H)$ its unitary group. For $U\in\mathrm U(H)$
  and $t\in(0,1]$ let $\kappa_t:L^2([0,t])\to L^2([0,1])$,
  $(\kappa_tv)(x)=t^{1/2}v(tx)$, and define $U_t=\kappa_t^{-1}U\kappa_t\oplus
  \Id_{L^2([t,1])}$, $U_0=\Id$. Show that each $U_t$ is unitary, that
  $\norm{U_tv-v}\le2\norm{v}_{L^2([0,t])}$, that $(UV)_t=U_tV_t$, and that
  $(U,t,v)\mapsto U_tv$ is continuous on $\mathrm U(H)\times[0,1]\times H$ when
  $\mathrm U(H)$ has the strong topology. Conclude that $\mathrm U(H)$ is
  contractible in the strong topology. (By Kuiper's theorem it is even
  contractible in the norm topology; this is not needed.)
\end{exercise}

\begin{exercise}[Smooth orthonormal frames]\label{scindex:ex:frame}
  Let $\phi:M\to B$ be a fibration of compact manifolds and $E\to M$ Hermitian.
  (a) Using an admissible embedding $J$ (\cref{scindex:lem:embedding}) and
  the finite rank approximation \ref{bott:B:finite} on $\Torus^N$, construct smooth
  families of finite rank orthogonal projections $\pi_k\in\Psido{-\infty}(M/B;E)$
  with $\norm{\pi_kA\pi_k-A}\to0$ uniformly on $B$ for every
  $A\in\Psido{-\infty}(M/B;E)$. (Hint: compress $J^*P_VJ$ and use the Riesz
  projection near $1$ of $J^*P_VJ$ for $V$ large.)
  (b) Alternatively, following Melrose, use \cref{scindex:ex:strong} and the fact
  that a fibre bundle with contractible fibre over a compact manifold has a
  section to find continuous sections $e'_i$ of the Hilbert bundle forming an
  orthonormal basis in each fibre; approximate them by smooth sections $e''_i$ with
  $\sup_b\norm{e'_i-e''_i}\le2^{-i-4}$, and show that the operator
  $T_bu=\sum_i\inner{u}{e'_i(b)}e''_i(b)$ satisfies $\norm{T_b-\Id}\le\frac14$, so
  that Gram--Schmidt applied to the $e''_i$ gives a smooth orthonormal basis.
\end{exercise}

\begin{exercise}\label{scindex:ex:relative}
  Let $P$, $Q$ be orthogonal projections on a Hilbert space with $P-Q$ compact.
  Show that $QP:\Ran P\to\Ran Q$ is Fredholm, and that for a pair of smooth
  families as in \cref{scindex:def:d} (with $B$ a point) its index equals
  $d(P,Q)$. (Reduce, by a homotopy and finite rank approximation, to the case
  $P=Q+\Pi$ with $\Pi$ of finite rank and $Q\Pi=0$.)
\end{exercise}

\begin{exercise}\label{scindex:ex:milnor-trace}
  In \cref{scindex:lem:milnor}, show that if $R_\pm$ are trace class then
  $\Tr(e_{A,B}-p_0)=\Tr(R_+^2)-\Tr(R_-^2)$, and that $\Tr(R_+^2)-\Tr(R_-^2)=
  \Tr(R_+)-\Tr(R_-)$ (use $AR_+=R_-A$ and $\Tr[A,BR_-]=0$ when the
  commutator is a commutator of a bounded and a trace class operator). Deduce
  \cref{scindex:prop:index-relative} for $B$ a point from the trace formula of
  \cref{index:thm:trace-formula}.
\end{exercise}

\begin{exercise}\label{scindex:ex:oscillatory}
  Let $A\in\Psisl{-\infty}(\R^n)$, $u\in\CIc(\R^n)$ and $f\in\CI(\R^n)$
  real-valued. Show that
  \begin{equation*}
    e^{-if/\epsilon}A(\epsilon)\bigl(e^{if/\epsilon}u\bigr)\longrightarrow
    \sigsl(A)(z,df(z))\,u(z)\quad\text{in }\CI(\R^n)\text{ as }\epsilon\downarrow0,
  \end{equation*}
  and use this to show that $\sigsl$ is invariant under changes of coordinates,
  in the sense of \ref{scindex:S-invariance}. Discuss what happens if
  $\epsilon$ is replaced by $2\epsilon$ (\cref{scindex:rem:slT}).
\end{exercise}

\begin{exercise}\label{scindex:ex:ground-state}
  Let $\mathcal V\to Z$ be a real vector bundle with a metric over a compact
  manifold, and $\gamma\in\cS(\mathcal V)$ the fibrewise normalized Gaussian.
  Show that $A\mapsto\Gamma A\Gamma^*$, $\Gamma u=\gamma\cdot\varpi^*u$, embeds
  $\Psido{-\infty}(Z;E)$ as a subalgebra of the operators on $\mathcal V$ with
  Schwartz kernels (Melrose's Proposition~12.3), that it maps $\Ginf(Z;E)$ into
  $\Id+\Psido{-\infty}_{\mathcal S}$, and that it preserves the classes in
  $K^{-1}(\mathrm{pt})$ of \cref{scindex:def:class}.
\end{exercise}

\begin{exercise}\label{scindex:ex:circle}
  For $M=\Torus$, $B=\mathrm{pt}$, compute $\Indsl{}$ of the class of the symbol of
  the operator $P_1=e^{-i\theta}\Pi_>+(\Id-\Pi_>)$ of \eqref{families:eq:P1} directly:
  show that a symbol representing $[\sigma_0(P_1)]$ can be chosen equal to $1$ for
  $\xi<-1$ and to $e^{-i\theta}$ for $\xi>1$, construct $e_{a,b}$, and show that
  the relative index of its quantization is $1$ by deforming it to the case
  $\epsilon=1$, where the quantization is a Fourier multiplier composed with
  multiplication operators.
\end{exercise}

\section*{Notes}
This chapter is based on Chapter~12 of R.~Melrose's notes \emph{From microlocal
to global analysis} (2008), which outlines a proof of the Atiyah--Singer index
theorem through semiclassical and adiabatic limits; many steps there are only
sketched, and several sections are empty. We have reorganized the material
around a definition of the even semiclassical index as a relative index of
idempotents, computed by the Toeplitz index of a loop
(\cref{scindex:def:d,scindex:def:Indsl}); with it, the comparison with the analytic
index (\cref{scindex:thm:analytic=sl}) reduces to Milnor's explicit form of the
index idempotent (\cref{scindex:lem:milnor}) and the multiplicativity of the
semiclassical symbol, in place of the combined `Toeplitz' calculus sketched in the
source. The passage from families of operators on the fibres to classes in
$K^{-1}(B)$ is made through an explicit isometric embedding of the fibres
(\cref{scindex:lem:embedding}) instead of Melrose's smooth orthonormal frames,
whose construction (via strong contractibility of the unitary group and the
existence of sections of bundles with contractible fibres) is left to
\cref{scindex:ex:strong,scindex:ex:frame}. The index of the Thom class
(\cref{scindex:thm:thom}) is computed from the harmonic oscillator, as suggested
by the source; this also completes the proof in \cref{ch:families}. The
multiplicativity of the semiclassical index under iterated fibrations, via the
adiabatic calculus, is only outlined (\cref{scindex:sec:adiabatic}); the final
proof uses instead the analytic multiplicativity of \cref{ch:families}. The
source's sections on Dirac families and spectral sections are empty and are
replaced by \cref{scindex:rem:dirac}.

The semiclassical approach to index theory has many antecedents; the use of the
harmonic oscillator to produce the Thom isomorphism goes back to the proof of Bott
periodicity by M.~F. Atiyah, \emph{Bott periodicity and the index of elliptic
operators}, Quart. J. Math. \textbf{19} (1968), 113--140, and the K-theoretic
setting is that of \cite{Atiyah-K,Atiyah-Singer-IV}. For the semiclassical
calculus see the references in \cref{ch:semiclassical}; for the isotropic calculus
and the harmonic oscillator, \cite{Shubin}. The idempotent \eqref{scindex:eq:milnor}
is the standard formula for the index map in algebraic K-theory, from J.~Milnor,
\emph{Introduction to algebraic K-theory}, Ann. of Math. Studies \textbf{72} (1971).
